\documentclass[11pt,a4paper]{article}

\usepackage[T1]{fontenc}
\usepackage{lmodern}
\usepackage[italian]{babel}
\usepackage{amsmath,amssymb,mathtools}
\usepackage{microtype}
\usepackage[a4paper,margin=2.6cm]{geometry}
\usepackage{setspace}
\usepackage{needspace}
\usepackage[most]{tcolorbox}
\usepackage{tikz}
\usetikzlibrary{arrows.meta, calc, backgrounds}

% Gestione interlinea (mantenere solo il comando desiderato)
\doublespacing

\setlength{\parindent}{1.2em}
\setlength{\parskip}{0.3em}
% Notazione
% Notazione
\newcommand{\dd}{\mathrm{d}}
\newcommand{\epz}{\varepsilon_0}
\newcommand{\Ev}{\vec{E}}
\newcommand{\Fv}{\vec{F}}
\newcommand{\uv}{\hat{u}}
\newcommand{\kc}{\ensuremath{\dfrac{1}{4\pi\epz}}}

% Titolo di paragrafo (maiuscoletto, senza grassetto né corsivo)
\newcommand{\titoletto}[1]{\par\Needspace{4\baselineskip}\medskip\noindent{\large\scshape #1}\par\nopagebreak\smallskip\noindent}

% Riquadro per i richiami (titolo senza grassetto)
\newtcolorbox{richiamo}[1]{
  enhanced, breakable,
  colback=gray!5, colframe=gray!55!black,
  boxrule=0.6pt, arc=2pt,
  left=8pt, right=8pt, top=6pt, bottom=6pt,
  fonttitle=\normalfont\large, coltitle=white,
  title={#1},
 before upper={\singlespacing }
}

\begin{document}

\begin{center}
{\LARGE Il campo elettrostatico nel vuoto\\[0.4ex]
dalla legge di Coulomb alle distribuzioni continue di carica}\\[1.2ex]
{\large Dispensa di Fisica 2 per la preparazione all'esame}
\end{center}

\vspace{1em}

\titoletto{Sintesi concettuale}
\noindent
Il percorso prende avvio dalla legge di Coulomb, che descrive la forza tra due cariche puntiformi ferme come azione diretta a distanza, e approda alla nozione di campo elettrostatico, che sposta l'attenzione dalla coppia di corpi alla regione di spazio che circonda la sorgente. Fissata la definizione operativa, il principio di sovrapposizione permette di costruire il campo di un insieme qualunque di cariche, tramite somme vettoriali per i sistemi discreti ed integrali di linea, di superficie e di volume per le distribuzioni continue. Seguono i casi canonici, cioè il dipolo, il filo indefinito, l'anello, il disco, il piano indefinito, il doppio strato e lo strato sferico, scelti perché in ciascuno la simmetria riduce il calcolo vettoriale a un solo integrale scalare. Una regolarità attraversa tutti gli esempi: a distanza grande rispetto alle dimensioni della sorgente, il campo coincide con quello di una carica puntiforme, mentre a distanza piccola emerge la geometria della distribuzione. La dispensa è scritta in modo da poter essere letta senza altri testi, i richiami di Meccanica e di Analisi sono raccolti negli appositi riquadri.

\medskip

\titoletto{Prerequisiti e strumenti di calcolo}
\begin{richiamo}{Richiamo Matematico/Fisica 1: forze, equilibrio, forza elastica ed energia}
Una forza è una grandezza vettoriale, misurata in newton, e la sua azione su un punto materiale di massa $m$ è regolata dalla relazione
\[
\Fv_{\mathrm{tot}}=m\vec a,
\]
dove $\Fv_{\mathrm{tot}}$ è la somma vettoriale di tutte le forze applicate e $\vec a$ l'accelerazione. Il terzo principio della dinamica afferma che, se un corpo A esercita su un corpo B la forza $\Fv_{AB}$, il corpo B esercita su A la forza
\[
\Fv_{BA}=-\Fv_{AB},
\]
con la stessa retta d'azione. Se $\Fv_{\mathrm{tot}}=\vec 0$, l'accelerazione è nulla, il punto resta fermo oppure procede di moto rettilineo uniforme, e in entrambi i casi la velocità è costante. Se invece $\Fv_{\mathrm{tot}}$ è costante e non nulla, l'accelerazione è costante e il moto, lungo una retta, risulta uniformemente accelerato
\[
x(t)=x_0+v_0t+\frac{1}{2}a t^2.
\]

La forza elastica
\[
F = -kx,
\]
con $k$ espresso in $\mathrm{N/m}$, è sempre diretta verso la posizione di equilibrio $x = 0$ e conduce all'equazione
\[
\ddot{x} = -\frac{k}{m}x,
\]
le cui soluzioni sono oscillazioni armoniche di pulsazione e periodo dati da
\[
\omega = \sqrt{\frac{k}{m}}, \qquad T = \frac{2\pi}{\omega}.
\]
Ogni volta che una forza risulta proporzionale e opposta allo spostamento, il moto descrive un oscillatore armonico, qualunque sia l'origine fisica della forza.

In contesti unidimensionali, le forze dipendenti esclusivamente dalla posizione risultano conservative, garantendo la conservazione dell'energia meccanica
\[
E_m = \frac{1}{2}mv^2 + U,
\]
con $U$ energia potenziale espressa in joule. L'energia potenziale si converte in cinetica e viceversa, ed il punto di inversione del moto si trova dove $v = 0$. In presenza di attrito parte dell'energia meccanica si dissipa, mentre in tutti gli esempi che seguono l'attrito è trascurato per dichiarazione esplicita.
\end{richiamo}

\begin{richiamo}{Richiamo Matematico/Fisica 1: vettori, versori, approssimazioni e geometria}
Un vettore $\vec r$ possiede modulo e direzione individuati da
\[
r=|\vec r|, \qquad \uv=\frac{\vec r}{r},
\]
con versore di modulo unitario e privo di dimensioni. La somma di vettori segue la regola del parallelogramma, o equivalentemente quella punta-coda, operativamente si scompone ogni vettore lungo gli assi scelti e si sommano le componenti omonime. Se tutti i vettori della somma, una volta riportati in catena punta-coda, formano un poligono chiuso, la risultante è nulla.

Per un triangolo qualsiasi di lati $a$, $b$, $c$ vale il teorema di Carnot
\[
c^2=a^2+b^2-2ab\cos\gamma,
\]
dove $\gamma$ è l'angolo opposto al lato $c$, nel caso $\gamma=\pi/2$ si ricade nel teorema di Pitagora. Per piccole quantità $\epsilon$ vale lo sviluppo di Taylor al primo ordine
\[
(1+\epsilon)^{\alpha}\simeq 1+\alpha\epsilon,
\]
con $\alpha$ reale, e tale approssimazione costituisce il principale strumento per ottenere il comportamento di un campo a grande distanza. Le aree utili risultano: circonferenza di raggio $s$, di lunghezza $2\pi s$; corona circolare di raggio interno $s$ e spessore infinitesimo $\dd s$, di area
\[
\dd S=2\pi s\,\dd s;
\]
superficie sferica di raggio $R$, di area $4\pi R^2$; fascia sferica compresa tra le colatitudini $\theta$ e $\theta+\dd\theta$, di area
\[
\dd S=2\pi R^2\sin\theta\,\dd\theta.
\]

Un integrale definito è il limite di una somma di contributi infinitesimi
\[
\int_a^b f(x)\,\dd x = F(b)-F(a).
\]
Quando una grandezza vettoriale è sommata su una distribuzione continua, l'integrale si scrive per componenti, e la simmetria della configurazione consente spesso di dimostrare in anticipo che alcune componenti si annullano a coppie. Gli angoli sono sempre espressi in radianti, un radiante è l'angolo che sottende un arco di lunghezza pari al raggio, per cui $\pi/2$ corrisponde a $90^{\circ}$ e $2\pi$ a $360^{\circ}$, e in Fisica 2 i gradi non compaiono nelle formule.
\end{richiamo}

\titoletto{La legge di Coulomb e il suo accertamento sperimentale}
La forza elettrica tra corpi carichi fermi è la prima interazione di cui si conosca una forma quantitativa esatta, ed il suo accertamento è sperimentale. Lo strumento è la bilancia di torsione, concettualmente affine al pendolo con cui Cavendish misurò l'attrazione gravitazionale tra masse: due sfere cariche, abbastanza piccole e distanti da poter essere schematizzate come punti, producono una torsione del filo proporzionale alla forza, e la ripetizione della misura a distanze e cariche differenti mostra che il modulo della forza è proporzionale al prodotto delle cariche e inversamente proporzionale al quadrato della distanza. La legge è dunque nata per cariche puntiformi, per corpi estesi si procede suddividendoli in elementi abbastanza piccoli da essere trattati come puntiformi, ed è questo il ruolo delle distribuzioni continue trattate più avanti.

\titoletto{Forma vettoriale della forza di Coulomb}
Siano $q$ e $q_0$ due cariche puntiformi, la prima posta in un punto $P$ e la seconda in un punto $P'$. Il vettore $\vec r$ va da $P$ a $P'$, ha modulo $r$ e individua il versore $\uv_r=\vec r/r$. La forza che $q$ esercita su $q_0$ vale
\[
\Fv=\frac{1}{4\pi\epz}\,\frac{q\,q'}{r^2}\,\uv_r ,
\qquad
\epz=8,8542\cdot 10^{-12}\,\mathrm{C^2/(N\,m^2)}.
\]

\begin{figure}[htbp]
  \centering
  \begin{tikzpicture}[>=Stealth, font=\small]

    % Stili personalizzati per elementi grafici
    \tikzset{
      gridline/.style={draw=gray!25, line width=0.3pt, dashed},
      badge/.style={draw=gray!40, fill=gray!10, rounded corners=3pt, inner xsep=8pt, inner ysep=4pt, font=\small\sffamily},
      forcearrow/.style={line width=1.3pt, ->},
      posarrow/.style={line width=0.8pt, ->, color=black!80}
    }

    % ==========================================
    % CASO 1: Repulsivo (q q' > 0)
    % ==========================================
    \begin{scope}[shift={(0,0)}]
      % Griglia tecnica di sfondo
      \draw[gridline] (-2.0,-1.1) grid (7.8,1.1);

      % Anelli concentrici (rappresentazione infografica di campo/potenziale)
      \foreach \r in {0.25, 0.45, 0.65} {
        \draw[gray!45, line width=0.4pt, dashed] (0,0) circle (\r);
        \draw[gray!45, line width=0.4pt, dashed] (4.2,0) circle (\r);
      }

      % Vettore Posizione r
      \draw[posarrow] (0,0) -- (4.2,0) node[midway, below=3pt] {$\vec{r}$};

      % Vettori Forza F' e F con prolungamento tratteggiato
      \draw[forcearrow, color=orange!90!red] (0,0) -- (-1.4,0) node[above=2pt] {$\vec{F}'$};
      \draw[dashed, color=orange!90!red, line width=0.8pt] (-1.4,0) -- (-1.8,0);

      \draw[forcearrow, color=orange!90!red] (4.2,0) -- (5.6,0) node[above=2pt] {$\vec{F}$};
      \draw[dashed, color=orange!90!red, line width=0.8pt] (5.6,0) -- (6.0,0);

      % Particelle / Cariche
      \filldraw[fill=orange!80!red, draw=black!80, line width=0.6pt] (0,0) circle (3.2pt);
      \filldraw[fill=purple!70!blue, draw=black!80, line width=0.6pt] (4.2,0) circle (3.2pt);

      % Etichette P, q, P', q'
      \node[above=4pt, font=\small\bfseries] at (0,0) {$P$};
      \node[below=4pt] at (0,0) {$q$};
      \node[above=4pt, font=\small\bfseries] at (4.2,0) {$P'$};
      \node[below=4pt] at (4.2,0) {$q'$};

      % Badge identificativo della condizione
      \node[badge] at (6.8,0) {$q\,q' > 0$};
    \end{scope}

    % ==========================================
    % CASO 2: Attrattivo (q q' < 0)
    % ==========================================
    \begin{scope}[shift={(0,-2.5)}]
      % Griglia tecnica di sfondo
      \draw[gridline] (-2.0,-1.1) grid (7.8,1.1);

      % Anelli concentrici
      \foreach \r in {0.25, 0.45, 0.65} {
        \draw[teal!40, line width=0.4pt, densely dotted] (0,0) circle (\r);
        \draw[teal!40, line width=0.4pt, densely dotted] (4.2,0) circle (\r);
      }

      % Vettore Posizione r
      \draw[posarrow] (0,0) -- (4.2,0) node[midway, below=3pt] {$\vec{r}$};

      % Vettori Forza F' e F (elevati per evitare sovrapposizioni)
      \draw[forcearrow, color=teal!80!black] (0,0.3) -- (1.4,0.3) node[midway, above=1pt] {$\vec{F}'$};
      \draw[dotted, color=teal!60, line width=0.6pt] (0,0) -- (0,0.3);

      \draw[forcearrow, color=teal!80!black] (4.2,0.3) -- (2.8,0.3) node[midway, above=1pt] {$\vec{F}$};
      \draw[dotted, color=teal!60, line width=0.6pt] (4.2,0) -- (4.2,0.3);

      % Particelle / Cariche
      \filldraw[fill=cyan!80!black, draw=black!80, line width=0.6pt] (0,0) circle (3.2pt);
      \filldraw[fill=blue!80!black, draw=black!80, line width=0.6pt] (4.2,0) circle (3.2pt);

      % Etichette P, q, P', q'
      \node[above=4pt, font=\small\bfseries] at (0,0) {$P$};
      \node[below=4pt] at (0,0) {$q$};
      \node[above=4pt, font=\small\bfseries] at (4.2,0) {$P'$};
      \node[below=4pt] at (4.2,0) {$q'$};

      % Badge identificativo della condizione
      \node[badge] at (6.8,0) {$q\,q' < 0$};
    \end{scope}

  \end{tikzpicture}
  \caption{Interazione elettrostatica tra due cariche puntiformi $q$ e $q'$ nei casi repulsivo ($q\,q' > 0$) ed attrattivo ($q\,q' < 0$).}
  \label{fig:legge_coulomb_infografica}
\end{figure}

Nel modulo compaiono $q$ e $q'$, i valori algebrici delle cariche in coulomb, e $r$, la distanza tra i punti in metri; $\epz$ è la costante dielettrica del vuoto, grandezza che caratterizza il mezzo e non le cariche; il fattore $1/(4\pi\epz)$, spesso indicato con $k$, vale circa $9\cdot 10^{9}\,\mathrm{N\,m^2/C^2}$, e la verifica dimensionale è immediata, perché $\mathrm{N\,m^2/C^2}\cdot\mathrm{C^2/m^2}=\mathrm{N}$. Il versore $\uv_r$ punta da $q$ verso $q_0$: se le cariche hanno lo stesso segno, il prodotto $qq_0$ è positivo e la forza su $q_0$ è concorde con $\uv_r$, quindi repulsiva; se hanno segno opposto, il prodotto è negativo e la forza è attrattiva. La forza che $q_0$ esercita su $q$ è uguale e opposta, come richiesto dal terzo principio, e giace sulla stessa retta. In aria la costante dielettrica differisce da $\epz$ in misura trascurabile, per cui la formula del vuoto vale anche per le applicazioni ordinarie.

\titoletto{Unità di misura della carica}
L'unità di carica è il coulomb, simbolo $\mathrm{C}$. Va ricordato che l'unità fondamentale del Sistema Internazionale è l'ampere, mentre il coulomb è un'unità derivata: un coulomb è la carica trasportata in un secondo da una corrente di un ampere, cioè $1\,\mathrm{C}=1\,\mathrm{A\,s}$. Il carattere derivato del coulomb è spesso rovesciato per errore, dato che l'ampere viene percepito come coulomb al secondo. Il valore assoluto della carica dell'elettrone, e dunque quella del protone, è $e=1,602\cdot 10^{-19}\,\mathrm{C}$, e conviene ricordarlo, perché compare in quasi ogni stima di ordine di grandezza. Per fissare le scale, si ricorda che una presa domestica italiana eroga nominalmente $230\,\mathrm{V}$ a $50\,\mathrm{Hz}$ (nelle lezioni è citato il valore storico di $220\,\mathrm{V}$) e che gli interruttori più comuni sopportano $16\,\mathrm{A}$.

\titoletto{Ordini di grandezza della forza coulombiana e legame ionico}
La forza di Coulomb è enorme rispetto alle scale della vita quotidiana. Due cariche di $1\,\mathrm{C}$ poste a $1\,\mathrm{m}$ di distanza si respingono con circa $9\cdot 10^{9}\,\mathrm{N}$, cioè con un peso equivalente a quello di quasi un milione di tonnellate: una situazione che i corpi macroscopici non sperimentano, perché sono globalmente neutri. Un'applicazione quotidiana si trova nel comportamento del sale da cucina in acqua. Il cloruro di sodio è un solido ionico, con ioni $\mathrm{Na}^+$ e $\mathrm{Cl}^-$ disposti in un reticolo cubico a facce centrate e separati da circa $0,28\,\mathrm{nm}$; la forza di attrazione tra una coppia di ioni, nel vuoto, vale
\[
F=\frac{1}{4\pi\epz}\frac{e^2}{r^2}\simeq 8,99\cdot 10^{9}\cdot\frac{(1,602\cdot 10^{-19})^2}{(0,28\cdot 10^{-9})^2}\ \mathrm{N}\simeq 2,9\cdot 10^{-9}\,\mathrm{N}.
\]
Tale valore, minuscolo in assoluto, è enorme alla scala del singolo ione, e il legame è per questo tenace. Immerso in acqua, la costante dielettrica $\epz$ va sostituita da $\epz\epsilon_r$ con $\epsilon_r\simeq 80$, e la forza tra gli ioni scende a circa $3,7\cdot 10^{-11}\,\mathrm{N}$. Il fenomeno ha anche una ragione geometrica: la molecola $\mathrm{H_2O}$ non è lineare, i due atomi di idrogeno formano un angolo con l'ossigeno, ed il baricentro della carica positiva non coincide con quello della negativa. La molecola è polare, si comporta da dipolo elettrico. I dipoli dell'acqua si orientano attorno agli ioni e schermano il legame, che si indebolisce fino a rompersi. Si distingue così il legame ionico dal legame covalente, nel quale gli elettroni sono messi in comune: nel diamante, dove ogni atomo di carbonio occupa il centro di un tetraedro, o nel quarzo e nei vetri silicatici, la rottura è molto più difficile, ed è questa la causa della durezza del diamante e delle proprietà isolanti del vetro.

\titoletto{Forza elettrica e forza gravitazionale a confronto}
Il confronto con la gravità, condotto sulla coppia di protoni, mostra quanto la forza elettrica prevalga a livello microscopico. In un nucleo due protoni distano dell'ordine di un femtometro, $1\,\mathrm{fm}=10^{-15}\,\mathrm{m}$. La repulsione elettrica vale
\[
F_e=\frac{1}{4\pi\epz}\frac{e^2}{r^2}\simeq 8,99\cdot 10^{9}\cdot\frac{(1,602\cdot 10^{-19})^2}{(10^{-15})^2}\ \mathrm{N}\simeq 2,3\cdot 10^{2}\,\mathrm{N},
\]
con un peso equivalente di circa $24\,\mathrm{kg}$. L'attrazione gravitazionale, calcolata con la costante di gravitazione universale $G=6,674\cdot 10^{-11}\,\mathrm{N\,m^2/kg^2}$ e con la massa del protone $m_p=1,673\cdot 10^{-27}\,\mathrm{kg}$, risulta
\[
F_g=G\,\frac{m_p^2}{r^2}\simeq 6,674\cdot 10^{-11}\cdot\frac{(1,673\cdot 10^{-27})^2}{(10^{-15})^2}\ \mathrm{N}\simeq 1,9\cdot 10^{-34}\,\mathrm{N},
\]
ed il rapporto $F_e/F_g$ è dell'ordine di $10^{36}$. Nel nucleo la gravità non agisce in modo apprezzabile, e la sola interazione elettrica condurrebbe alla dispersione immediata dei protoni. Poiché la materia esiste, si conclude che nel nucleo agiscono altre interazioni, la forte e la debole; i neutroni, privi di carica e quindi indifferenti alla repulsione di Coulomb, contribuiscono al legame. L'equilibrio è però delicato: al crescere del numero atomico servono sempre più neutroni, e gli elementi più pesanti risultano instabili. L'uranio, attraverso una lunga catena che passa per il radon, termina nel piombo, che è stabile. L'elettrone non contribuisce in modo sostanziale a questo equilibrio, trovandosi lontano dal nucleo per una distanza dell'ordine di $10^4$--$10^5$ volte le dimensioni nucleari.

\titoletto{La gravità tra corpi ordinari e celesti}
Con la stessa procedura si apprezza il peso della gravità per corpi ordinari. Due persone di massa $70\,\mathrm{kg}$ poste a $1\,\mathrm{m}$ si attraggono con
\[
F=G\,\frac{m^2}{r^2}=6,674\cdot 10^{-11}\cdot 4900\ \mathrm{N}\simeq 3,3\cdot 10^{-7}\,\mathrm{N},
\]
equivalenti al peso di una massa di circa $33\,\mathrm{\mu g}$. La stessa interazione diventa rilevante solo con masse planetarie, il peso di una persona di $70\,\mathrm{kg}$ sulla Terra è di circa $7\cdot 10^{2}\,\mathrm{N}$, attrazione dovuta al pianeta. L'abitudine a eseguire simili stime numeriche rientra tra le competenze richieste in sede d'esame.

\titoletto{Impostazione dei problemi e principio di sovrapposizione delle forze}
Per l'impostazione dei problemi conviene accompagnare ogni domanda con un disegno: si collocano la carica $q$ in $P$ e la carica $q_0$ in $P'$, si traccia la distanza $r$, si sceglie il versore $\uv_r$ e si rappresenta $\Fv_{q\to q_0}$ applicata in $P'$ nel verso appropriato. Il principio di sovrapposizione stabilisce che, dato un insieme di cariche $q_i$ fisse, la forza risultante su $q_0$ è la somma vettoriale delle forze esercitate da ciascuna
\[
\Fv=\sum_i \Fv_i=\sum_i \frac{1}{4\pi\epz}\frac{q_i\,q_0}{r_i^2}\,\uv_i=q_0\,\sum_i \frac{1}{4\pi\epz}\frac{q_i}{r_i^2}\,\uv_i ,
\]
dove $r_i$ e $\uv_i$ sono distanza e versore riferiti alla carica $q_i$. Il fattore $q_0$ si porta fuori dalla somma, definendo l'origine del campo. Come esempio, disponendo quattro cariche uguali sui vertici di un quadrato ed una quinta nel centro, ciascuna delle prime quattro dista dal centro di mezza diagonale, per cui le quattro forze hanno uguale modulo; a due a due, quelle dei vertici opposti sono uguali e contrarie, e la risultante è nulla. La carica centrale, soggetta a $\Fv_{\mathrm{tot}}=\vec 0$, resta ferma oppure procede di moto rettilineo uniforme.

\begin{figure}[htbp]
  \centering
  \begin{tikzpicture}[>=Stealth, font=\small\sffamily]

    % Stili grafici
    \tikzset{
      gridline/.style={draw=gray!20, line width=0.3pt, dashed},
      badge/.style={draw=gray!40, fill=gray!10, rounded corners=3pt, inner xsep=8pt, inner ysep=4pt, font=\small\sffamily},
      forcearrow/.style={line width=1.3pt, ->},
      posline/.style={draw=gray!50, line width=0.7pt, dashed},
      unitarrow/.style={line width=1.0pt, ->, color=black!85}
    }

    % Griglia tecnica di sottofondo
    \draw[gridline] (-4.6,-2.4) grid (4.6,4.6);

    % Coordinate dei centri di carica
    \coordinate (q0) at (0,0);
    \coordinate (q1) at (3.4, 1.2);
    \coordinate (q2) at (0.0, 3.8);
    \coordinate (qi) at (-2.6, 3.2);
    \coordinate (qn) at (-3.6, 0.8);

    % Anelli concentrici attorno alle cariche
    \foreach \r in {0.25, 0.45, 0.65} {
      \draw[orange!50!red, line width=0.4pt, dashed] (q0) circle (\r);
      \draw[orange!50!red, line width=0.4pt, dashed] (q1) circle (\r);
      \draw[orange!50!red, line width=0.4pt, dashed] (q2) circle (\r);
      \draw[orange!50!red, line width=0.4pt, dashed] (qi) circle (\r);
      \draw[teal!50, line width=0.4pt, densely dotted] (qn) circle (\r);
    }

    % Congiungenti tratteggiate
    \draw[posline] (q1) -- (q0);
    \draw[posline] (q2) -- (q0);
    \draw[posline] (qi) -- (q0);
    \draw[posline] (qn) -- (q0);

    % Distanze r_k
    \node[below right=1pt, font=\small\itshape, color=gray!80!black] at ($(q1)!0.55!(q0)$) {$r_1$};
    \node[right=3pt, font=\small\itshape, color=gray!80!black] at ($(q2)!0.55!(q0)$) {$r_2$};
    \node[above right=1pt, font=\small\itshape, color=gray!80!black] at ($(qi)!0.55!(q0)$) {$r_i$};
    \node[below=2pt, font=\small\itshape, color=gray!80!black] at ($(qn)!0.55!(q0)$) {$r_n$};

    % Versori u_k lungo le congiungenti
    \draw[unitarrow] ($(q1)!0cm!(q0)$) -- ($(q1)!1cm!(q0)$) node[midway, below=0 pt, font=\bfseries\small] {$\vec{u}_1$};
    \draw[unitarrow] ($(q2)!0cm!(q0)$) -- ($(q2)!1cm!(q0)$) node[midway, right=3pt, font=\bfseries\small] {$\vec{u}_2$};
    \draw[unitarrow] ($(qi)!0cm!(q0)$) -- ($(qi)!1cm!(q0)$) node[midway, above right=-1pt, font=\bfseries\small] {$\vec{u}_i$};
    \draw[unitarrow] ($(qn)!0cm!(q0)$) -- ($(qn)!1cm!(q0)$) node[midway, above=1pt, font=\bfseries\small] {$\vec{u}_n$};

    % Vettori forza applicati alla carica sonda q0
    \draw[forcearrow, color=orange!90!red] (q0) -- ($(q0)!-2.0cm!(q1)$) node[below left=-1pt, font=\bfseries] {$\vec{F}_1$};
    \draw[forcearrow, color=orange!90!red] (q0) -- ($(q0)!-2.0cm!(q2)$) node[below=2pt, font=\bfseries] {$\vec{F}_2$};
    \draw[forcearrow, color=orange!90!red] (q0) -- ($(q0)!-2.0cm!(qi)$) node[below right=-1pt, font=\bfseries] {$\vec{F}_i$};
    \draw[forcearrow, color=teal!80!black] (q0) -- ($(q0)!2.0cm!(qn)$) node[above left=-1pt, font=\bfseries] {$\vec{F}_n$};

    % Particelle cariche
    \filldraw[fill=orange!80!red, draw=black!80, line width=0.6pt] (q1) circle (3.5pt);
    \node[above right=3pt, font=\small\bfseries] at (q1) {$q_1$};

    \filldraw[fill=orange!80!red, draw=black!80, line width=0.6pt] (q2) circle (3.5pt);
    \node[right=4pt, font=\small\bfseries] at (q2) {$q_2$};

    \filldraw[fill=orange!80!red, draw=black!80, line width=0.6pt] (qi) circle (3.5pt);
    \node[above left=3pt, font=\small\bfseries] at (qi) {$q_i$};

    \filldraw[fill=teal!70!black, draw=black!80, line width=0.6pt] (qn) circle (3.5pt);
    \node[below left=3pt, font=\small\bfseries] at (qn) {$q_n$};

    \filldraw[fill=purple!70!blue, draw=black!80, line width=0.6pt] (q0) circle (4.0pt);
    \node[below right=3pt, font=\small\bfseries] at (q0) {$q_0$};

    % Badge superiore esplicativo
    \node[badge] at (0, 5.2) {$\vec{F}_{\mathrm{tot}} = \sum_{k} \vec{F}_k$};

  \end{tikzpicture}
  \caption{Principio di sovrapposizione per le forze elettrostatiche agenti sulla carica di prova $q_0$.}
  \label{fig:sovrapposizione_infografica}
\end{figure}


\titoletto{Definizione del campo elettrostatico}
Il campo elettrostatico nasce dall'osservazione che la parte dipendente dalle sorgenti è separabile dalla carica che subisce la forza. Il campo $\Ev$ generato in un punto da un sistema di cariche ferme è definito come la forza risultante su una carica di prova positiva $q_0$ divisa per $q_0$, con il limite
\[
\Ev=\lim_{q_0\to 0}\frac{\Fv}{q_0}.
\]



Il campo è un vettore, misurato in $\mathrm{N/C}$ (equivalente a $\mathrm{V/m}$), e rappresenta la forza per unità di carica positiva. Il limite $q_0\to 0$ evita che la sonda perturbi la distribuzione primaria. Si distingue così il campo elettrostatico, generato da cariche ferme, dal campo elettrico in senso generale.

\titoletto{Portata e limiti del concetto di campo}
%\section{Portata e limiti del concetto di campo elettrico}

La definizione del campo elettrico poggia su due principi concettuali fondamentali:
\begin{enumerate}
	\item \textbf{Indipendenza dalle sorgenti:} Il campo è svincolato dalle modalità specifiche della sua generazione; tale proprietà rende il modello idoneo a descrivere non solo le configurazioni elettrostatiche, ma anche i campi indotti e le cariche in movimento.
	\item \textbf{Realtà fisica oggettiva:} Oltre all'aspetto operativo legato alla misura delle forze su cariche sonda, il campo costituisce un'entità fisica autonoma. La sua esistenza prescinde dall'effettiva presenza di una carica di prova nella regione interessata.
\end{enumerate}

%\subsection*{Limiti della definizione operativa}
Nel regime elettrostatico, la rigorosa definizione del vettore intensità di campo $\vec{E}$ richiede che sia le dimensioni spaziali sia la carica $q$ del corpo di prova siano infinitesimali rispetto alle sorgenti. Tuttavia, la natura quantizzata della materia impone un limite fisico invalicabile: la carica non può essere inferiore al quanto fondamentale $e$ (carica dell'elettrone).

Nei sistemi microscopici, caratterizzati da poche cariche elementari, il passaggio al limite perde di significato operativo. La formulazione classica del campo tramite carica sonda mantiene perciò piena validità esclusivamente nei contesti macroscopici, dove la carica totale è il risultato dell'aggregazione di un numero estremamente elevato di unità discrete.

%\subsection*{Scomposizione dell'interazione}
\titoletto{Significato fisico dell'intensità di campo}
L'intensità del campo elettrico in un punto rappresenta la forza esercitata sull'unità di carica positiva. La forza $\vec{F}$ agente su una carica puntiforme $q$ collocata in tale punto è espressa da:
\begin{equation}
	\vec{F} = q\vec{E}
\end{equation}
sotto l'ipotesi che la presenza della carica $q$ non alteri la distribuzione spaziale delle sorgenti circostanti.

\titoletto{Azione a distanza e campo come mediatore dell'interazione}
Mentre la legge di Coulomb descrive l'interazione mediante un'azione a distanza immediata tra cariche:
\[
\text{sorgente} \longrightarrow \text{forza su un corpo carico}
\]
il concetto di campo scinde il fenomeno in due fasi distinte e fisicamente autonome:
\[
\text{sorgente} \longrightarrow \text{campo elettrico} \longrightarrow \text{forza sulla carica } q
\]

Questa bipartizione offre due prospettive analitiche complementari:
\begin{itemize}
	\item \textbf{Determinazione dalle sorgenti:} Nota la distribuzione delle cariche sorgente, il campo elettrico può essere calcolato in ogni punto senza alcun riferimento a cariche di prova.
	\item \textbf{Predizione della dinamica:} Conoscendo il campo in una regione, è possibile prevedere la forza esercitata su una qualsiasi carica di prova, prescindendo dalla conoscenza delle sorgenti che lo hanno generato.
\end{itemize}

%\subsection*{Dall'Elettrostatica all'Elettrodinamica}
Nei contesti stazionari (elettrostatica), la descrizione coulombiana a distanza e quella basata sul campo conducono a risultati identici; in questo ambito, il campo potrebbe apparire come un puro artificio matematico.

\titoletto{Realtà fisica del campo e velocità finita dei segnali}
La necessità fisica del campo come entità reale ed autonoma emerge in modo irrefutabile nei fenomeni dinamici. Se una carica sorgente viene accelerata, la teoria dell'azione a distanza prevederebbe una variazione istantanea della forza percepita da una seconda carica distante. Sperimentalmente si osserva invece un ritardo di propagazione: l'accelerazione degli elettroni in un'antenna trasmittente agisce sugli elettroni di un'antenna ricevente posta a distanza $d$ solo dopo un intervallo temporale:
\begin{equation}
	\tau = \frac{d}{c}
\end{equation}
dove $c$ rappresenta la velocità della luce nel vuoto. Il campo elettrico si rivela quindi un'entità fisica reale, dotata di energia ed impulso, capace di propagare le perturbazioni nello spazio a velocità finita.
La ragione per cui si preferisce il campo alla forza è la velocità finita di propagazione dei segnali, limitata da $c\simeq 3\cdot 10^{8}\,\mathrm{m/s}$. La descrizione avviene in due stadi: la sorgente crea il campo nello spazio, e il campo agisce sul corpo carico con un ritardo temporale. Un vantaggio pratico consiste nel determinare il campo indipendentemente da una specifica carica sonda.

\titoletto{Campo di una carica puntiforme}
Data una  sorgente puntiforme $q$ ed un punto $P$ a distanza $r$, si ha
\[
\Ev(P)=\frac{1}{4\pi\epz}\,\frac{q}{r^2}\,\uv_r ,
\]
dove $q$ è la carica sorgente in coulomb, $r$ la distanza in metri e $\uv_r$ il versore radiale. Per $q>0$ il campo è uscente, per $q<0$ è entrante. Il modulo vale $E=|q|/(4\pi\epz r^2)$, decrescendo come $r^{-2}$.

\begin{figure}[htbp]
	\centering
	\begin{tikzpicture}[>=Stealth, font=\small\sffamily]
		
		% Stili grafici personalizzati
		\tikzset{
			gridline/.style={draw=gray!20, line width=0.3pt, dashed},
			badge/.style={draw=gray!40, fill=gray!10, rounded corners=3pt, inner xsep=8pt, inner ysep=4pt, font=\small\sffamily},
			fieldarrow/.style={line width=1.3pt, ->},
			posline/.style={draw=gray!50, line width=0.7pt, dashed},
			unitarrow/.style={line width=1.0pt, ->, color=black!85}
		}
		
		% Griglia tecnica di sfondo
		\draw[gridline] (-1.8,-2.2) grid (6.8,4.6);
		
		% ==========================================
		% CASO 1: Carica positiva (q > 0) - Campo uscente
		% ==========================================
		\coordinate (A1) at (0, 2.2);
		\coordinate (B1) at (4.0, 3.7);
		
		% Anelli concentrici di campo
		\foreach \r in {0.25, 0.45, 0.65} {
			\draw[orange!50!red, line width=0.4pt, dashed] (A1) circle (\r);
		}
		
		% Linea della distanza r
		\draw[posline] (A1) -- (B1) node[midway, below right=1pt, font=\small\itshape, color=gray!80!black] {$r$};
		
		% Versore u
		\draw[unitarrow] ($(A1)!0cm!(B1)$) -- ($(A1)!1cm!(B1)$) node[midway, below right=-1pt, font=\bfseries\small] {$\vec{u}$};
		
		% Punto P e vettore Campo Elettrico E (uscente)
		\filldraw[fill=black!70, draw=black] (B1) circle (1.5pt) node[above left=2pt, font=\small\itshape] {$P$};
		\draw[fieldarrow, color=orange!90!red] (B1) -- ($(B1)!-1.6cm!(A1)$) node[above right=-1pt, font=\bfseries] {$\vec{E}$};
		
		% Carica sorgente q
		\filldraw[fill=orange!80!red, draw=black!80, line width=0.6pt] (A1) circle (3.5pt);
		\node[above left=3pt, font=\small\bfseries] at (A1) {$q$};
		
		% Badge identificativo della condizione
		\node[badge] at (5.8, 2.2) {$q > 0$};
		
		
		% ==========================================
		% CASO 2: Carica negativa (q1 < 0) - Campo entrante
		% ==========================================
		\coordinate (A2) at (0, -1.0);
		\coordinate (B2) at (4.0, 0.5);
		
		% Anelli concentrici di campo
		\foreach \r in {0.25, 0.45, 0.65} {
			\draw[teal!50, line width=0.4pt, densely dotted] (A2) circle (\r);
		}
		
		% Linea della distanza r
		\draw[posline] (A2) -- (B2) node[midway, below right=1pt, font=\small\itshape, color=gray!80!black] {$r$};
		
		% Versore u
		\draw[unitarrow] ($(A2)!0cm!(B2)$) -- ($(A2)!1cm!(B2)$) node[midway, below right=-1pt, font=\bfseries\small] {$\vec{u}$};
		
		% Punto P' e vettore Campo Elettrico E (entrante)
		\filldraw[fill=black!70, draw=black] (B2) circle (1.5pt) node[above left=2pt, font=\small\itshape] {$P'$};
		\draw[fieldarrow, color=teal!80!black] (B2) -- ($(B2)!1.6cm!(A2)$) node[below right=-1pt, font=\bfseries] {$\vec{E}$};
		
		% Carica sorgente q1
		\filldraw[fill=teal!70!black, draw=black!80, line width=0.6pt] (A2) circle (3.5pt);
		\node[above left=3pt, font=\small\bfseries] at (A2) {$q_1$};
		
		% Badge identificativo della condizione
		\node[badge] at (5.8, -1.0) {$q_1 < 0$};
		
	\end{tikzpicture}
	\caption{Orientazione del campo elettrostatico $\vec{E}$ nel punto di misura per una carica sorgente positiva $q$ ed una carica negativa $q_1$.}
	\label{fig:campo_carica_puntiforme_infografica}
\end{figure}


\titoletto{Campo di un sistema di cariche puntiformi}
Il principio di sovrapposizione per un sistema di $N$ cariche puntiformi fornisce
\[
\Ev=\sum_{i=1}^{N}\Ev_i=\sum_{i=1}^{N}\frac{1}{4\pi\epz}\,\frac{q_i}{r_i^2}\,\uv_i .
\]
\titoletto{Esempio: quattro cariche uguali ai vertici di un quadrato}
Nei vertici di un quadrato con quattro cariche uguali, il campo nel centro risulta nullo per cancellazione a coppie. 

\begin{figure}[htbp]
  \centering
  \begin{tikzpicture}[>=Stealth, font=\small\sffamily]

    % Stili grafici
    \tikzset{
      gridline/.style={draw=gray!20, line width=0.3pt, dashed},
      badge/.style={draw=gray!40, fill=gray!10, rounded corners=3pt, inner xsep=8pt, inner ysep=4pt, font=\small\sffamily},
      fieldarrow/.style={line width=1.3pt, ->},
      posline/.style={draw=gray!40, line width=0.6pt, dashed},
      geomline/.style={draw=gray!60, line width=0.8pt}
    }

    % Griglia tecnica di sfondo
    \draw[gridline] (-3.8,-3.2) grid (3.8,3.6);

    % Coordinate vertici (quadrato) e centro
    \coordinate (C)  at (0,0);
    \coordinate (q1) at (-2.3, 2.3);
    \coordinate (q2) at ( 2.3, 2.3);
    \coordinate (q3) at ( 2.3,-2.3);
    \coordinate (q4) at (-2.3,-2.3);

    % Perimetro del quadrato
    \draw[geomline] (q1) -- (q2) -- (q3) -- (q4) -- cycle;

    % Anelli concentrici attorno alle cariche
    \foreach \q in {q1, q2, q3, q4} {
      \foreach \r in {0.25, 0.45, 0.65} {
        \draw[orange!50!red, line width=0.4pt, dashed] (\q) circle (\r);
      }
    }

    % Diagonali / Linee d'azione
    \draw[posline] (q1) -- (q3);
    \draw[posline] (q2) -- (q4);

    % Vettori campo elettrico al centro (uscenti dalle cariche positive)
    % E1 da q1 punta verso q3
    \draw[fieldarrow, color=orange!90!red] (C) -- ($(C)!1.5cm!(q3)$) node[below right=-2pt, font=\bfseries] {$\vec{E}_1$};
    % E2 da q2 punta verso q4
    \draw[fieldarrow, color=orange!90!red] (C) -- ($(C)!1.5cm!(q4)$) node[below left=-2pt, font=\bfseries] {$\vec{E}_2$};
    % E3 da q3 punta verso q1
    \draw[fieldarrow, color=orange!90!red] (C) -- ($(C)!1.5cm!(q1)$) node[above left=-2pt, font=\bfseries] {$\vec{E}_3$};
    % E4 da q4 punta verso q2
    \draw[fieldarrow, color=orange!90!red] (C) -- ($(C)!1.5cm!(q2)$) node[above right=-2pt, font=\bfseries] {$\vec{E}_4$};

    % Punto centrale P
    \filldraw[fill=black!70, draw=black] (C) circle (1.8pt) node[below=6pt, font=\small\itshape] {$P$};

    % Cariche sorgente nei vertici
    \filldraw[fill=orange!80!red, draw=black!80, line width=0.6pt] (q1) circle (3.5pt) node[above left=2pt, font=\small\bfseries] {$q_1$};
    \filldraw[fill=orange!80!red, draw=black!80, line width=0.6pt] (q2) circle (3.5pt) node[above right=2pt, font=\small\bfseries] {$q_2$};
    \filldraw[fill=orange!80!red, draw=black!80, line width=0.6pt] (q3) circle (3.5pt) node[below right=2pt, font=\small\bfseries] {$q_3$};
    \filldraw[fill=orange!80!red, draw=black!80, line width=0.6pt] (q4) circle (3.5pt) node[below left=2pt, font=\small\bfseries] {$q_4$};

    % Badge esplicativo
    \node[badge] at (0, 3.1) {$\vec{E}_{\mathrm{tot}} = \vec{E}_1 + \vec{E}_2 + \vec{E}_3 + \vec{E}_4 = \vec{0}$};

  \end{tikzpicture}
  \caption{Cancellazione a coppie dei campi elettrostatici al centro di una configurazione quadrata di quattro cariche uguali.}
  \label{fig:campo_quadrato_infografica}
\end{figure}

\titoletto{Esempio: tre cariche uguali ai vertici di un triangolo equilatero}
Al centro di un triangolo equilatero con tre cariche uguali, i tre campi formano angoli di $120^{\circ}=2\pi/3$. Fissando l'asse $y$ lungo uno dei tre vettori, gli altri due hanno componenti $y$ pari a $E\cos(2\pi/3)=-E/2$, annullando la componente totale.

\begin{figure}[htbp]
  \centering
  \begin{tikzpicture}[>=Stealth, font=\small\sffamily]

    % Stili grafici
    \tikzset{
      gridline/.style={draw=gray!20, line width=0.3pt, dashed},
      badge/.style={draw=gray!40, fill=gray!10, rounded corners=3pt, inner xsep=8pt, inner ysep=4pt, font=\small\sffamily},
      fieldarrow/.style={line width=1.3pt, ->},
      posline/.style={draw=gray!40, line width=0.6pt, dashed},
      geomline/.style={draw=gray!60, line width=0.8pt},
      axis/.style={draw=gray!70, line width=0.8pt, ->}
    }

    % Griglia tecnica di sfondo
    \draw[gridline] (-3.8,-3.4) grid (3.8,3.4);

    % Coordinate vertici (triangolo equilatero, R = 2.5 cm) e centro C
    \coordinate (C)  at (0,0);
    \coordinate (q1) at (0, -2.5);                   % 270 gradi
    \coordinate (q2) at (2.165, 1.25);               % 30 gradi
    \coordinate (q3) at (-2.165, 1.25);              % 150 gradi

    % Perimetro del triangolo
    \draw[geomline] (q1) -- (q2) -- (q3) -- cycle;

    % Assi cartesiani orientati (asse y lungo E1)
    \draw[axis] (-3.0, 0) -- (3.2, 0) node[right, font=\small\itshape] {$x$};
    \draw[axis] (0, -3.0) -- (0, 2.8) node[above, font=\small\itshape] {$y$};

    % Anelli concentrici attorno alle cariche
    \foreach \q in {q1, q2, q3} {
      \foreach \r in {0.25, 0.45, 0.65} {
        \draw[orange!50!red, line width=0.4pt, dashed] (\q) circle (\r);
      }
    }

    % Congiungenti dai vertici al centro
    \draw[posline] (q1) -- (C);
    \draw[posline] (q2) -- (C);
    \draw[posline] (q3) -- (C);

    % Vettori Campo Elettrico al centro (lunghezza 1.8 cm)
    % E1 diretto lungo +y (allontana da q1)
    \draw[fieldarrow, color=orange!90!red] (C) -- (0, 1.8) node[right=2pt, font=\bfseries] {$\vec{E}_1$};
    % E2 diretto verso il 3° quadrante (210°)
    \draw[fieldarrow, color=orange!90!red] (C) -- (-1.5588, -0.9) node[below left=-1pt, font=\bfseries] {$\vec{E}_2$};
    % E3 diretto verso il 4° quadrante (330°)
    \draw[fieldarrow, color=orange!90!red] (C) -- (1.5588, -0.9) node[below right=-1pt, font=\bfseries] {$\vec{E}_3$};

    % Proiezioni tratteggiate delle componenti y di E2 ed E3 sull'asse y
    \draw[gray!60, line width=0.6pt, densely dotted] (-1.5588, -0.9) -- (0, -0.9);
    \draw[gray!60, line width=0.6pt, densely dotted] (1.5588, -0.9) -- (0, -0.9);
    \filldraw[fill=orange!90!red, draw=black] (0, -0.9) circle (1.2pt);
    \node[left=3pt, font=\scriptsize\sffamily, color=gray!90!black] at (0, -0.9) {$-E/2$};

    % Arco per l'angolo tra E1 ed E2 (2\pi/3)
    \draw[line width=0.5pt, color=gray!80!black] (0, 0.7) arc (90:210:0.7);
    \node[font=\scriptsize\sffamily, color=gray!90!black] at (-0.6, 0.6) {$2\pi/3$};

    % Cariche sorgente nei vertici
    \filldraw[fill=orange!80!red, draw=black!80, line width=0.6pt] (q1) circle (3.5pt) node[below=4pt, font=\small\bfseries] {$q_1$};
    \filldraw[fill=orange!80!red, draw=black!80, line width=0.6pt] (q2) circle (3.5pt) node[above right=2pt, font=\small\bfseries] {$q_2$};
    \filldraw[fill=orange!80!red, draw=black!80, line width=0.6pt] (q3) circle (3.5pt) node[above left=2pt, font=\small\bfseries] {$q_3$};

    % Punto centrale P
    \filldraw[fill=black!70, draw=black] (C) circle (1.8pt) node[above right=2pt, font=\small\itshape] {$P$};

    % Badge identificativo della componente y
    \node[badge] at (0, -3.8) {$E_y = E - \frac{E}{2} - \frac{E}{2} = 0$};

  \end{tikzpicture}
  \caption{Simmetria dei tre campi elettrostatici a $2\pi/3$ e scomposizione lungo l'asse $y$ al centro di un triangolo equilatero.}
  \label{fig:campo_triangolo_infografica}
\end{figure}

\titoletto{Distribuzioni continue di carica}
Per una distribuzione continua l'integrale vettoriale assume la forma
\[
\Ev=\frac{1}{4\pi\epz}\int \frac{\dd q}{r^2}\,\uv_r ,
\]
con elementi di carica espressi da
\[
\dd q=\lambda\,\dd l, \qquad \dd q=\sigma\,\dd S, \qquad \dd q=\rho\,\dd\tau,
\]
rispettivamente per densità lineare $\lambda$ ($\mathrm{C/m}$),

\begin{figure}[htbp]
  \centering
  \begin{tikzpicture}[>=Stealth, font=\small\sffamily]

    % Stili grafici
    \tikzset{
      gridline/.style={draw=gray!20, line width=0.3pt, dashed},
      badge/.style={draw=gray!40, fill=gray!10, rounded corners=3pt, inner xsep=8pt, inner ysep=4pt, font=\small\sffamily},
      fieldarrow/.style={line width=1.3pt, ->, color=orange!90!red},
      posline/.style={draw=gray!50, line width=0.7pt, dashed},
      curveline/.style={draw=orange!80!black, line width=1.5pt}
    }

    % Griglia tecnica di sfondo
    \draw[gridline] (-2.2,-1.6) grid (6.8,3.8);

    % Coordinate chiave: P' lungo la curva e punto P nello spazio
    \coordinate (P1) at (-1.5, -0.5);
    \coordinate (P2) at (0.0, 1.2);    % Punto sorgente P'
    \coordinate (P3) at (2.0, 2.5);
    \coordinate (P)  at (4.8, 0.8);    % Punto di misura P

    % Curva rappresentante il filo carico l
    \draw[curveline] (P1) .. controls (-0.8, 0.2) and (-0.5, 0.8) .. (P2)
                         .. controls (0.6, 1.7) and (1.2, 2.3) .. (P3);

    % Anelli concentrici di campo attorno all'elemento sorgente dl in P'
    \foreach \r in {0.25, 0.45, 0.65} {
      \draw[orange!50!red, line width=0.4pt, dashed] (P2) circle (\r);
    }

    % Elemento infinitesimo dl evidenziato sulla curva
    \draw[line width=3.2pt, color=orange!90!red] (0.0, 1.2) -- (0.22, 1.4);

    % Congiungente r dal punto sorgente P' al punto P
    \draw[posline] (P2) -- (P) node[midway, above right=-1pt, font=\small\itshape, color=gray!80!black] {$\vec{r}$};

    % Vettore campo infinitesimo dE generato in P
    \filldraw[fill=black!70, draw=black] (P) circle (1.8pt) node[below right=2pt, font=\small\itshape] {$P$};
    \draw[fieldarrow] (P) -- ($(P)!-1.8cm!(P2)$) node[above right=-1pt, font=\bfseries] {$\mathrm{d}\vec{E}$};

    % Etichette dell'elemento sorgente P' e della curva l
    \filldraw[fill=orange!80!red, draw=black!80, line width=0.6pt] (P2) circle (3.0pt);
    \node[above left=3pt, font=\small\bfseries] at (P2) {$P'$};
    \node[above left=1pt, font=\small\bfseries, color=orange!90!red] at (0.22, 1.4) {$\mathrm{d}l$};
    \node[right=2pt, font=\small\bfseries, color=orange!80!black] at (P3) {$l$};

    % Badge della carica infinitesima
    \node[badge] at (5.2, 3.1) {$\mathrm{d}q = \lambda\,\mathrm{d}l$};

  \end{tikzpicture}
  \caption{Elemento di carica infinitesimo $\mathrm{d}q = \lambda\,\mathrm{d}l$ su una distribuzione lineare $l$ e campo $\mathrm{d}\vec{E}$ generato nel punto $P$.}
  \label{fig:distribuzione_lineare_infografica}
\end{figure}

 superficiale $\sigma$ ($\mathrm{C/m^2}$)

\begin{figure}[htbp]
  \centering
  \begin{tikzpicture}[>=Stealth, font=\small\sffamily]

    % Stili grafici
    \tikzset{
      gridline/.style={draw=gray!20, line width=0.3pt, dashed},
      badge/.style={draw=gray!40, fill=gray!10, rounded corners=3pt, inner xsep=8pt, inner ysep=4pt, font=\small\sffamily},
      fieldarrow/.style={line width=1.3pt, ->, color=orange!90!red},
      posline/.style={draw=gray!50, line width=0.7pt, dashed}
    }

    % Griglia tecnica di sfondo
    \draw[gridline] (-2.2,-2.2) grid (6.8,3.8);

    % Punti geometrici della superficie curva
    \coordinate (S1) at (-1.8, 1.2);
    \coordinate (S2) at (1.5, 3.0);
    \coordinate (S3) at (3.8, -1.2);
    \coordinate (S4) at (0.5, -1.8);

    % Disegno della superficie con sfumatura calda e contorno definito
    \filldraw[fill=orange!12, draw=orange!80!black, line width=1.1pt]
      (S1) .. controls (-0.2, 2.2) and (0.8, 2.6) .. (S2)
           .. controls (2.6, 1.2) and (3.2, 0.2) .. (S3)
           .. controls (2.1, -1.6) and (1.2, -1.6) .. (S4)
           .. controls (-0.8, -1.0) and (-1.4, 0.0) .. (S1);

    % Coordinate sorgente P' e punto di misura P
    \coordinate (P2) at (0.4, 0.6);   % Punto sorgente P'
    \coordinate (P)  at (5.0, 1.4);   % Punto P

    % Anelli concentrici di campo attorno a P'
    \foreach \r in {0.25, 0.45, 0.65} {
      \draw[orange!50!red, line width=0.4pt, dashed] (P2) circle (\r);
    }

    % Patch dell'elemento infinitesimo dS attorno a P'
    \filldraw[fill=orange!80!red, draw=black!80, line width=0.6pt]
      ($(P2)+(-0.35,-0.2)$) -- ($(P2)+(0.25,-0.25)$) -- ($(P2)+(0.38,0.22)$) -- ($(P2)+(-0.22,0.27)$) -- cycle;

    % Congiungente r da P' a P
    \draw[posline] (P2) -- (P) node[midway, above right=-1pt, font=\small\itshape, color=gray!80!black] {$\vec{r}$};

    % Vettore campo dE applicato in P
    \filldraw[fill=black!70, draw=black] (P) circle (1.8pt) node[below right=2pt, font=\small\itshape] {$P$};
    \draw[fieldarrow] (P) -- ($(P)!-1.8cm!(P2)$) node[above right=-1pt, font=\bfseries] {$\mathrm{d}\vec{E}$};

    % Etichette P', dS ed S
    \filldraw[fill=black!80] (P2) circle (1.8pt) node[below left=2pt, font=\small\bfseries] {$P'$};
    \node[above=4pt, font=\small\bfseries, color=orange!90!red] at (P2) {$\mathrm{d}S$};
    \node[below left=2pt, font=\small\bfseries, color=orange!80!black] at (S3) {$S$};

    % Badge della carica infinitesima
    \node[badge] at (5.2, 3.1) {$\mathrm{d}q = \sigma\,\mathrm{d}S$};

  \end{tikzpicture}
  \caption{Elemento di superficie $\mathrm{d}S$ con carica $\mathrm{d}q = \sigma\,\mathrm{d}S$ appartenente alla distribuzione bidimensionale $S$.}
  \label{fig:distribuzione_superficiale_infografica}
\end{figure}

 e volumica $\rho$ ($\mathrm{C/m^3}$).

\begin{figure}[htbp]
  \centering
  \begin{tikzpicture}[>=Stealth, font=\small\sffamily]

    % Stili grafici
    \tikzset{
      gridline/.style={draw=gray!20, line width=0.3pt, dashed},
      badge/.style={draw=gray!40, fill=gray!10, rounded corners=3pt, inner xsep=8pt, inner ysep=4pt, font=\small\sffamily},
      fieldarrow/.style={line width=1.3pt, ->, color=orange!90!red},
      posline/.style={draw=gray!50, line width=0.7pt, dashed}
    }

    % Griglia tecnica di sfondo
    \draw[gridline] (-2.2,-2.4) grid (6.8,4.2);

    % Contorno del volume 3D \tau
    \coordinate (V1) at (-0.8, 3.2);
    \coordinate (V2) at (1.8, 3.6);
    \coordinate (V3) at (2.2, 0.2);
    \coordinate (V4) at (0.2, -1.8);
    \coordinate (V5) at (-1.8, -0.6);

    % Riempimento sfumato ed equipotenziale del solido 3D
    \filldraw[fill=orange!10, draw=orange!80!black, line width=1.1pt]
      (V1) .. controls (0.4, 3.8) and (1.2, 3.8) .. (V2)
           .. controls (2.6, 2.2) and (2.8, 1.2) .. (V3)
           .. controls (1.8, -1.0) and (1.2, -1.6) .. (V4)
           .. controls (-0.8, -2.2) and (-1.4, -1.4) .. (V5)
           .. controls (-2.2, 0.6) and (-1.6, 2.2) .. (V1);

    % Tratto di prospettiva interna per tridimensionalità
    \draw[gray!60, line width=0.7pt, dashed]
      (-1.7, 0.8) .. controls (-0.2, 0.2) and (0.8, 0.4) .. (2.4, 1.2);

    % Coordinate sorgente P' e punto di misura P
    \coordinate (P2) at (0.2, 2.0);   % Centro del cubetto d\tau
    \coordinate (P)  at (5.0, 0.8);   % Punto P

    % Anelli concentrici di campo attorno all'elemento volumetrico
    \foreach \r in {0.28, 0.50, 0.72} {
      \draw[orange!50!red, line width=0.4pt, dashed] (P2) circle (\r);
    }

    % Cubetto volumetrico 3D (d\tau) in assonometria
    \begin{scope}[shift={(P2)}]
      % Facce posteriori / interne
      \filldraw[fill=orange!30, draw=orange!90!red, line width=0.5pt] (-0.2,-0.2) rectangle (0.2,0.2);
      % Faccia superiore
      \filldraw[fill=orange!50, draw=orange!90!red, line width=0.5pt] (-0.2,0.2) -- (-0.05,0.38) -- (0.35,0.38) -- (0.2,0.2) -- cycle;
      % Faccia laterale destra
      \filldraw[fill=orange!70!red, draw=orange!90!red, line width=0.5pt] (0.2,-0.2) -- (0.35,-0.02) -- (0.35,0.38) -- (0.2,0.2) -- cycle;
    \end{scope}

    % Congiungente r dal punto sorgente P' al punto P
    \draw[posline] (P2) -- (P) node[midway, above right=-1pt, font=\small\itshape, color=gray!80!black] {$\vec{r}$};

    % Vettore campo dE applicato in P
    \filldraw[fill=black!70, draw=black] (P) circle (1.8pt) node[below right=2pt, font=\small\itshape] {$P$};
    \draw[fieldarrow] (P) -- ($(P)!-1.8cm!(P2)$) node[above right=-1pt, font=\bfseries] {$\mathrm{d}\vec{E}$};

    % Etichette P', d\tau ed il volume complessivo \tau
    \filldraw[fill=black!80] (P2) circle (1.6pt) node[below left=2pt, font=\small\bfseries] {$P'$};
    \node[above=6pt, font=\small\bfseries, color=orange!90!red] at (P2) {$\mathrm{d}\tau$};
    \node[below left=2pt, font=\small\bfseries, color=orange!80!black] at (V4) {$\tau$};

    % Badge della carica infinitesima
    \node[badge] at (5.2, 3.3) {$\mathrm{d}q = \rho\,\mathrm{d}\tau$};

  \end{tikzpicture}
  \caption{Elemento volumetrico $\mathrm{d}\tau$ con carica infinitesima $\mathrm{d}q = \rho\,\mathrm{d}\tau$ inserito all'interno del volume triplo $\tau$.}
  \label{fig:distribuzione_volumetrica_infografica}
\end{figure}


\begin{richiamo}{Richiamo Matematico: sviluppi in serie per distanze grandi e angoli piccoli}
Nel confronto con una sorgente puntiforme si ricorre a due approssimazioni. Per piccoli angoli $|\theta|\ll 1$ in radianti si ha
\[
\sin\theta\simeq\theta, \qquad \cos\theta\simeq 1-\frac{\theta^2}{2}.
\]
Lo sviluppo binomiale per $|\epsilon|\ll 1$ fornisce
\[
(1+\epsilon)^{\alpha}\simeq 1+\alpha\epsilon,
\]
da cui seguono ad esempio le relazioni
\[
(1\pm\epsilon)^{-2}\simeq 1\mp 2\epsilon, \qquad (1+\epsilon)^{-1/2}\simeq 1-\frac{\epsilon}{2}.
\]
\end{richiamo}

\titoletto{Il dipolo elettrico}
Per il principio di sovrapposizione, il campo $\vec{E}_{o}$ prodotto da un dipolo in un punto qualsiasi $P$ è la somma dei contributi delle due cariche che lo compongono:
\begin{equation}
	\vec{E}_{o}=\vec{E}_{o+}+\vec{E}_{o-}.
\end{equation}
Qui $\vec{E}_{o+}$ è il campo della carica $+q$ ed $\vec{E}_{o-}$ quello della carica $-q$, poste a distanza $l$ l'una dall'altra e con momento di dipolo $\vec{p}=q\,\vec{l}$, diretto dalla carica negativa a quella positiva.

\begin{figure}[htbp]
	\centering
	\begin{tikzpicture}[>=Stealth, font=\small\sffamily]
		
		% Stili grafici
		\tikzset{
			gridline/.style={draw=gray!20, line width=0.3pt, dashed},
			badge/.style={draw=gray!40, fill=gray!10, rounded corners=3pt, inner xsep=8pt, inner ysep=4pt, font=\small\sffamily},
			fieldarrow/.style={line width=1.3pt, ->},
			posline/.style={draw=gray!50, line width=0.7pt, dashed},
			vectorp/.style={line width=1.5pt, ->, color=black!85}
		}
		
		% Griglia tecnica di sfondo
		\draw[gridline] (-3.8,-1.6) grid (5.8,5.2);
		
		% Coordinate cariche e punto P
		\coordinate (A) at (-2.5, 0);  % -q
		\coordinate (B) at ( 1.5, 0);  % +q
		\coordinate (P) at ( 2.6, 3.6); % Punto P
		
		% Segmento dipolo e momento p
		\draw[line width=0.8pt, color=gray!70] (A) -- (B);
		\draw[vectorp] (-1.2, 0) -- (0.5, 0) node[midway, above=2pt, font=\small\bfseries] {$\vec{p}$};
		
		% Anelli concentrici attorno alle cariche
		\foreach \r in {0.25, 0.45, 0.65} {
			\draw[teal!50, line width=0.4pt, densely dotted] (A) circle (\r);
			\draw[orange!50!red, line width=0.4pt, dashed] (B) circle (\r);
		}
		
		% Connessioni tratteggiate verso P
		\draw[posline] (A) -- (P);
		\draw[posline] (B) -- (P);
		
		% Etichetta distanza l
		\node[below=4pt, font=\small\itshape, color=gray!80!black] at ($(A)!0.5!(B)$) {$l$};
		
		% Vettori campo in P
		% E_0+ (uscente da +q, diretto lungo B -> P)
		\draw[fieldarrow, color=orange!90!red] (P) -- ($(P)!-1.8cm!(B)$) node[above right=-1pt, font=\bfseries] {$\vec{E}_{0+}$};
		% E_0- (entrante in -q, diretto lungo P -> A)
		\draw[fieldarrow, color=teal!80!black] (P) -- ($(P)!1.5cm!(A)$) node[below left=-1pt, font=\bfseries] {$\vec{E}_{0-}$};
		
		% Punto P
		\filldraw[fill=black!70, draw=black] (P) circle (1.8pt) node[right=4pt, font=\small\itshape] {$P$};
		
		% Cariche sorgente
		\filldraw[fill=teal!70!black, draw=black!80, line width=0.6pt] (A) circle (3.5pt) node[below left=2pt, font=\small\bfseries] {$-q$};
		\filldraw[fill=orange!80!red, draw=black!80, line width=0.6pt] (B) circle (3.5pt) node[below right=2pt, font=\small\bfseries] {$+q$};
		
		% Badge descrittivo
		\node[badge] at (4.0, 4.6) {$\vec{E}_0(P) = \vec{E}_{0+} + \vec{E}_{0-}$};
		
	\end{tikzpicture}
	\caption{Campo elettrostatico generato da un dipolo elettrico in un punto generico $P$ dello spazio.}
	\label{fig:dipolo_generico_infografica}
\end{figure}

Conviene studiare due situazioni particolari, dette prima e seconda posizione principale di Gauss. Nella prima il punto $P$ giace sulla retta che passa per le due cariche (l'asse del dipolo). In tal caso $\vec{E}_{o+}$ ed $\vec{E}_{o-}$ hanno la stessa direzione di quella retta e versi opposti.
Detta $r$ la distanza di $P$ dal centro $O$ del dipolo, il modulo del campo risultante è
\begin{equation}
	E_{o1}=\frac{1}{4\pi\epsilon_{0}}\left[\frac{q}{\left(r-\frac{l}{2}\right)^{2}}-\frac{q}{\left(r+\frac{l}{2}\right)^{2}}\right]=\frac{1}{4\pi\epsilon_{0}}\,\frac{2\,r\,l\,q}{\left(r^{2}-\frac{l^{2}}{4}\right)^{2}}.
\end{equation}
Il campo ha il verso di $\vec{p}$, perché il contributo della carica più vicina, cioè $+q$, prevale su quello dell'altra. Per punti lontani dal dipolo, $l\ll r$, si possono trascurare i termini di ordine superiore in $l/r$ e si ottiene
\begin{equation}
	\vec{E}_{o1}=\frac{1}{4\pi\epsilon_{0}}\,\frac{2\,\vec{p}}{r^{3}}.
\end{equation}






\begin{figure}[htbp]
	\centering
	\begin{tikzpicture}[>=Stealth, font=\small\sffamily]
		
		% Stili grafici
		\tikzset{
			gridline/.style={draw=gray!20, line width=0.3pt, dashed},
			badge/.style={draw=gray!40, fill=gray!10, rounded corners=3pt, inner xsep=8pt, inner ysep=4pt, font=\small\sffamily},
			fieldarrow/.style={line width=1.3pt, ->},
			posline/.style={draw=gray!50, line width=0.7pt, dashed},
			vectorp/.style={line width=1.5pt, ->, color=black!85},
			dimline/.style={line width=0.6pt, <->, color=gray!80!black}
		}
		
		% Griglia tecnica di sfondo
		\draw[gridline] (-3.8,-1.8) grid (6.8,2.2);
		
		% Coordinate cariche e punto P
		\coordinate (A) at (-2.8, 0);  % -q
		\coordinate (B) at (-0.8, 0);  % +q
		\coordinate (C) at (-1.8, 0);  % Centro dipolo
		\coordinate (P) at ( 3.8, 0);  % Punto P sull'asse
		
		% Asse e prolungamento
		\draw[line width=0.8pt, color=gray!70] (A) -- (B);
		\draw[posline] (B) -- (P);
		
		% Momento di dipolo p
		\draw[vectorp] (-2.3, 0) -- (-1.3, 0) node[midway, above=2pt, font=\small\bfseries] {$\vec{p}$};
		
		% Anelli concentrici
		\foreach \r in {0.25, 0.45, 0.65} {
			\draw[teal!50, line width=0.4pt, densely dotted] (A) circle (\r);
			\draw[orange!50!red, line width=0.4pt, dashed] (B) circle (\r);
		}
		
		% Vettori campo in P (asse x)
		% E_0+ a destra (uscente da +q, piu intenso)
		\draw[fieldarrow, color=orange!90!red] (P) -- ($(P)+(1.8,0)$) node[above=1pt, font=\bfseries] {$\vec{E}_{0+}$};
		% E_0- a sinistra (entrante in -q, meno intenso)
		\draw[fieldarrow, color=teal!80!black] (P) -- ($(P)+(-1.1,0)$) node[above=1pt, font=\bfseries] {$\vec{E}_{0-}$};
		
		% Punto P
		\filldraw[fill=black!70, draw=black] (P) circle (1.8pt) node[above right=3pt, font=\small\itshape] {$P$};
		
		% Cariche
		\filldraw[fill=teal!70!black, draw=black!80, line width=0.6pt] (A) circle (3.5pt) node[above left=2pt, font=\small\bfseries] {$-q$};
		\filldraw[fill=orange!80!red, draw=black!80, line width=0.6pt] (B) circle (3.5pt) node[above right=2pt, font=\small\bfseries] {$+q$};
		
		% Linea di quota distanza r
		\draw[dimline] (-1.8, -1.0) -- (3.8, -1.0) node[midway, fill=white, inner sep=2pt, font=\small\itshape] {$r$};
		\draw[gray!60, line width=0.5pt] (-1.8, -0.3) -- (-1.8, -1.2);
		\draw[gray!60, line width=0.5pt] (3.8, -0.3) -- (3.8, -1.2);
		
		% Badge descrittivo
		\node[badge] at (1.5, 1.5) {$E_{01} \simeq \dfrac{2p}{4\pi\varepsilon_0 r^3} \quad (r \gg l)$};
		
	\end{tikzpicture}
	\caption{Campo elettrostatico del dipolo nei punti appartenenti all'asse longitudinale.}
	\label{fig:dipolo_asse_infografica}
\end{figure}

Nella seconda posizione il punto $P$ si trova sull'asse geometrico del dipolo, cioè sulla retta perpendicolare al segmento che congiunge le cariche e passante per il suo centro. In questo caso:
\begin{enumerate}
	\item[1)] i due campi hanno lo stesso modulo,
	\begin{equation}
		E_{o+}=E_{o-}=\frac{1}{4\pi\epsilon_{0}}\,\frac{q}{r^{2}+\frac{l^{2}}{4}};
	\end{equation}
	\item[2)] le componenti perpendicolari alla direzione del dipolo sono uguali e opposte, quindi si annullano;
	\item[3)] le componenti parallele alla direzione del dipolo sono uguali e concordi, quindi si sommano.
\end{enumerate}



mentre sul piano mediano a distanza $y\gg a$ si riduce a
\[
E_{02}=-\frac{q\,l}{4\pi\epz\,(r^2+l^2/4)^{3/2}}\simeq-\frac{p}{4\pi\epz\,r^3}.
\]
In entrambe le direzioni il campo decresce come $r^{-3}$.

\begin{figure}[htbp]
	\centering
	\begin{tikzpicture}[>=Stealth, font=\small\sffamily]
		
		% Stili grafici
		\tikzset{
			gridline/.style={draw=gray!20, line width=0.3pt, dashed},
			badge/.style={draw=gray!40, fill=gray!10, rounded corners=3pt, inner xsep=8pt, inner ysep=4pt, font=\small\sffamily},
			fieldarrow/.style={line width=1.3pt, ->},
			posline/.style={draw=gray!50, line width=0.7pt, dashed},
			vectorp/.style={line width=1.5pt, ->, color=black!85},
			axisline/.style={draw=gray!70, line width=0.8pt, dashed}
		}
		
		% Griglia tecnica di sfondo
		\draw[gridline] (-3.8,-1.2) grid (3.8,6.2);
		
		% Coordinate cariche e punto P
		\coordinate (A) at (-2.0, 0);  % -q
		\coordinate (B) at ( 2.0, 0);  % +q
		\coordinate (C) at ( 0.0, 0);  % Centro del dipolo
		\coordinate (P) at ( 0.0, 4.2); % Punto P sul piano mediano
		
		% Asse mediano verticale
		\draw[axisline] (C) -- (0, 5.4);
		
		% Segmento dipolo e momento p
		\draw[line width=0.8pt, color=gray!70] (A) -- (B);
		\draw[vectorp] (-1.2, 0) -- (1.2, 0) node[midway, below=2pt, font=\small\bfseries] {$\vec{p}$};
		
		% Anelli concentrici
		\foreach \r in {0.25, 0.45, 0.65} {
			\draw[teal!50, line width=0.4pt, densely dotted] (A) circle (\r);
			\draw[orange!50!red, line width=0.4pt, dashed] (B) circle (\r);
		}
		
		% Connessioni tratteggiate verso P
		\draw[posline] (A) -- (P);
		\draw[posline] (B) -- (P);
		
		% Etichetta quota r
		\node[left=2pt, font=\small\itshape, color=gray!80!black] at (0, 2.1) {$r$};
		
		% Definizione vettori campo in P
		\coordinate (Eplus)  at ($(P)!-1.8cm!(B)$);
		\coordinate (Eminus) at ($(P)!1.8cm!(A)$);
		\coordinate (Eres)   at ($(Eplus)+(Eminus)-(P)$);
		
		% Parallelogramma delle componenti
		\draw[gray!50, line width=0.6pt, dashed] (Eplus) -- (Eres) -- (Eminus);
		\draw[gray!40, line width=0.5pt, densely dotted] (Eplus) -- ($(P)!(Eplus)!(0,5.4)$);
		\draw[gray!40, line width=0.5pt, densely dotted] (Eminus) -- ($(P)!(Eminus)!(0,5.4)$);
		
		% Vettori campo
		\draw[fieldarrow, color=orange!90!red] (P) -- (Eplus) node[above left=-1pt, font=\bfseries] {$\vec{E}_{0+}$};
		\draw[fieldarrow, color=teal!80!black] (P) -- (Eminus) node[below left=-1pt, font=\bfseries] {$\vec{E}_{0-}$};
		\draw[fieldarrow, color=purple!80!black] (P) -- (Eres) node[left=3pt, font=\bfseries] {$\vec{E}_0$};
		
		% Punto P
		\filldraw[fill=black!70, draw=black] (P) circle (1.8pt) node[right=4pt, font=\small\itshape] {$P$};
		
		% Cariche sorgente
		\filldraw[fill=teal!70!black, draw=black!80, line width=0.6pt] (A) circle (3.5pt) node[below left=2pt, font=\small\bfseries] {$-q$};
		\filldraw[fill=orange!80!red, draw=black!80, line width=0.6pt] (B) circle (3.5pt) node[below right=2pt, font=\small\bfseries] {$+q$};
		
		% Badge descrittivo
		\node[badge] at (3, 5.0) {$E_{02} \simeq -\dfrac{p}{4\pi\varepsilon_0 r^3} \quad (r \gg l)$};
		
	\end{tikzpicture}
	\caption{Composizione dei vettori di campo elettrostatico sul piano mediano del dipolo con risultante $\vec{E}_0$ antiparallela a $\vec{p}$.}
	\label{fig:dipolo_mediano_infografica}
\end{figure}



\titoletto{Filo rettilineo uniforme indefinito}
Si consideri una retta $x$ di lunghezza illimitata, sulla quale la carica è distribuita con densità lineare costante
\begin{equation}
	\lambda=\frac{dq}{dx}=\text{cost.}
\end{equation}
Si vuole determinare il campo in un punto $P$ che dista $a$ dalla retta; $O$ indica il piede della perpendicolare condotta da $P$ alla retta stessa.

\begin{figure}[htbp]
	\centering
	\begin{tikzpicture}[>=Stealth, font=\small\sffamily]
		
		% Stili grafici
		\tikzset{
			gridline/.style={draw=gray!20, line width=0.3pt, dashed},
			badge/.style={draw=gray!40, fill=gray!10, rounded corners=3pt, inner xsep=8pt, inner ysep=4pt, font=\small\sffamily},
			fieldarrow/.style={line width=1.3pt, ->},
			posline/.style={draw=gray!50, line width=0.7pt, dashed},
			axisline/.style={draw=gray!70, line width=0.8pt},
			curveline/.style={draw=orange!80!black, line width=1.5pt},
			dimline/.style={line width=0.6pt, <->, color=gray!80!black}
		}
		
		% Griglia tecnica di sfondo
		\draw[gridline] (-2.4,-4.6) grid (5.6,4.6);
		
		% Coordinate principali
		\coordinate (O)          at (0, 0);
		\coordinate (P)          at (3.2, 0);       % Punto P a distanza radiale a
		\coordinate (Pprime)     at (0, 2.4);       % Centro dell'elemento dy
		\coordinate (Pprime_bot) at (0, 2.15);      % Estremo inferiore di dy
		\coordinate (Pprime_top) at (0, 2.65);      % Estremo superiore di dy
		\coordinate (Pdouble)    at (0, -2.4);      % Elemento simmetrico -y
		
		% Filo rettilineo verticale (asse y)
		\draw[curveline] (0, -4.2) -- (0, 4.2) node[above=2pt, font=\small\bfseries, color=orange!80!black] {$y$};
		
		% Anelli concentrici di campo attorno agli elementi dy e dy'
		\foreach \r in {0.25, 0.45, 0.65} {
			\draw[orange!50!red, line width=0.4pt, dashed] (Pprime) circle (\r);
			\draw[teal!50, line width=0.4pt, densely dotted] (Pdouble) circle (\r);
		}
		
		% Evidenziazione dell'elemento infinitesimo dy
		\draw[line width=3.5pt, color=orange!90!red] (Pprime_bot) -- (Pprime_top);
		\draw[line width=3.5pt, color=teal!70!black] (0, -2.65) -- (0, -2.15);
		
		% Congiungenti dal punto P agli estremi dell'elemento dy per definire d\theta
		\draw[posline, color=orange!80!red!60] (P) -- (Pprime_bot);
		\draw[posline, color=orange!80!red!60] (P) -- (Pprime_top);
		\draw[posline] (Pprime) -- (P);
		\draw[posline] (Pdouble) -- (P);
		\draw[axisline, dashed] (O) -- (P);
		
		% Etichette delle distanze r ed a
		\node[above right=-1pt, font=\small\itshape, color=gray!80!black] at ($(Pprime)!0.55!(P)$) {$r$};
		\node[below=2pt, font=\small\itshape, color=gray!80!black] at ($(O)!0.5!(P)$) {$a$};
		
		% Quota della coordinata y
		\draw[dimline] (-0.7, 0) -- (-0.7, 2.4) node[midway, fill=white, inner sep=1.5pt, font=\small\itshape] {$y$};
		\draw[gray!60, line width=0.5pt] (-0.2, 0) -- (-0.9, 0);
		\draw[gray!60, line width=0.5pt] (-0.2, 2.4) -- (-0.9, 2.4);
		
		% Vettori campo infinitesimo applicati in P
		\draw[fieldarrow, color=orange!90!red] (P) -- ($(P)!-1.8cm!(Pprime)$) node[below right=-1pt, font=\bfseries] {$\mathrm{d}\vec{E}_0$};
		\draw[fieldarrow, color=teal!80!black] (P) -- ($(P)!-1.8cm!(Pdouble)$) node[above right=-1pt, font=\bfseries] {$\mathrm{d}\vec{E}_0'$};
		
		% Componente risultante perpendicolare dE_perp
		\coordinate (E1) at ($(P)!-1.8cm!(Pprime)$);
		\coordinate (E2) at ($(P)!-1.8cm!(Pdouble)$);
		\coordinate (Eres) at ($(E1)+(E2)-(P)$);
		\draw[gray!40, line width=0.6pt, dashed] (E1) -- (Eres) -- (E2);
		\draw[fieldarrow, color=purple!80!black] (P) -- (Eres) node[right=2pt, font=\bfseries] {$\mathrm{d}\vec{E}_\perp$};
		
		% Angolo theta principale nel vertice P
		\draw[line width=0.5pt, color=gray!80!black] (P) +(-0.8,0) arc (180:143.13:0.8);
		\node[font=\scriptsize\sffamily, color=gray!90!black] at ($(P)+(-1.05,0.20)$) {$\theta$};
		
		% Angolo infinitesimo d\theta sotteso dall'elemento dy nel punto P
		\draw[line width=0.7pt, color=orange!90!red] (P) +({180-atan2(2.15,3.2)}:1.65cm) arc ({180-atan2(2.15,3.2)}:{180-atan2(2.65,3.2)}:1.65cm);
		\node[font=\scriptsize\sffamily\bfseries, color=orange!90!red] at ($(P)+(-1.95, 1.28)$) {$\mathrm{d}\theta$};
		
		% Angoli theta per la decomposizione dei vettori di campo
		\draw[line width=0.5pt, color=gray!80!black] (P) +(0.8,0) arc (0:-36.87:0.8);
		\node[font=\scriptsize\sffamily, color=gray!90!black] at ($(P)+(1.05,-0.20)$) {$\theta$};
		
		\draw[line width=0.5pt, color=gray!80!black] (P) +(0.8,0) arc (0:36.87:0.8);
		\node[font=\scriptsize\sffamily, color=gray!90!black] at ($(P)+(1.05,0.20)$) {$\theta$};
		
		% Punti ed etichette sul filo
		\filldraw[fill=black!70, draw=black] (O) circle (1.8pt) node[below left=2pt, font=\small\bfseries] {$O$};
		\filldraw[fill=black!70, draw=black] (P) circle (1.8pt) node[above=4pt, font=\small\itshape] {$P$};
		
		\filldraw[fill=orange!80!red, draw=black!80, line width=0.6pt] (Pprime) circle (3.0pt);
		\node[right=4pt, font=\small\bfseries, color=orange!90!red] at (Pprime) {$\mathrm{d}y$};
		\node[left=2pt, font=\small\bfseries] at (Pprime) {$P'$};
		
		\filldraw[fill=teal!70!black, draw=black!80, line width=0.6pt] (Pdouble) circle (3.0pt);
		\node[right=4pt, font=\small\bfseries, color=teal!80!black] at (Pdouble) {$\mathrm{d}y'$};
		
		% Badge con la relazione differenziale tra coordinata ed angolo
		\node[badge] at (2.6, 3.9) {$\mathrm{d}y = \dfrac{a\,\mathrm{d}\theta}{\cos^2\theta}$};
		
	\end{tikzpicture}
	\caption{Rappresentazione dell'angolo infinitesimo $\mathrm{d}\theta$ sotteso dall'elemento $\mathrm{d}y$ del filo verticale nel punto $P$.}
	\label{fig:campo_filo_verticale_dtheta_infografica}
\end{figure}

\begin{figure}[htbp]
	\centering
	\begin{tikzpicture}[>=Stealth, font=\small\sffamily]
		
		% Stili grafici
		\tikzset{
			gridline/.style={draw=gray!20, line width=0.3pt, dashed},
			badge/.style={draw=gray!40, fill=gray!10, rounded corners=3pt, inner xsep=8pt, inner ysep=4pt, font=\small\sffamily},
			posline/.style={draw=gray!50, line width=0.7pt, dashed},
			curveline/.style={draw=orange!80!black, line width=1.5pt},
			zoombox/.style={draw=orange!60!black, line width=0.8pt, fill=orange!5, rounded corners=6pt}
		}
		
		% Griglia tecnica di sfondo
		\draw[gridline] (-5.4,-3.2) grid (5.4,3.8);
		
		% ==========================================
		% SINISTRA: Schema contestuale con lente
		% ==========================================
		\begin{scope}[shift={(-3.6, 0)}, scale=0.6]
			% Filo verticale
			\draw[curveline] (0, -4.2) -- (0, 4.2) node[above=1pt, font=\scriptsize\bfseries, color=orange!80!black] {$y$};
			\coordinate (Pmini) at (3.2, 0);
			\coordinate (Pprimemini) at (0, 2.2);
			
			% Congiungenti di contesto
			\draw[posline] (0,0) -- (Pmini) node[midway, below, font=\tiny] {$a$};
			\draw[posline] (Pprimemini) -- (Pmini);
			
			% Elemento dy minimizzato
			\draw[line width=2.5pt, color=orange!90!red] (0, 1.9) -- (0, 2.5);
			\filldraw[fill=orange!80!red] (Pprimemini) circle (2.5pt);
			\node[left=2pt, font=\tiny\bfseries, color=orange!90!red] at (Pprimemini) {$\mathrm{d}y$};
			
			% Cerchio lente d'ingrandimento
			\draw[draw=orange!90!red, line width=1.0pt, fill=white, fill opacity=0.15] (Pprimemini) circle (1.0cm);
			
			% Linee di proiezione verso il riquadro di ingrandimento
			\draw[dashed, color=orange!60!black, line width=0.5pt] (0.6, 2.9) -- (2.0, 3.2);
			\draw[dashed, color=orange!60!black, line width=0.5pt] (0.6, 1.5) -- (2.0, -2.8);
		\end{scope}
		
		% ==========================================
		% DESTRA: Ingrandimento del triangolo rettangolo
		% ==========================================
		\begin{scope}[shift={(0.8, 0)}]
			% Riquadro di ingrandimento
			\filldraw[zoombox] (-1.6, -2.8) rectangle (4.2, 3.2);
			
			% Coordinate del triangolo ingrandito (\theta = 35 gradi)
			\coordinate (A) at (0, 2.0);                  % Estremo superiore di dy
			\coordinate (B) at (0, -2.0);                 % Estremo inferiore di dy
			\coordinate (C) at (1.879, 0.683);            % Piede della perpendicolare su r
			
			% Segmento dy (Ipotenusa)
			\draw[line width=3.5pt, color=orange!90!red] (A) -- (B);
			\filldraw[fill=orange!80!red, draw=black!80, line width=0.5pt] (A) circle (2.5pt);
			\filldraw[fill=orange!80!red, draw=black!80, line width=0.5pt] (B) circle (2.5pt);
			
			% Raggi diretti verso P (a grande distanza verso basso-destra, -35 gradi)
			\draw[line width=1.1pt, color=black!75, ->] (A) -- ($(A) + (-35:3.8cm)$) node[right=2pt, font=\small\itshape] {verso $P$};
			\draw[line width=1.1pt, color=black!75, ->] (B) -- ($(B) + (-35:3.8cm)$);
			
			% Perpendicolare dal punto inferiore B al raggio superiore (Cateto r d\theta)
			\draw[line width=1.4pt, color=purple!80!black] (B) -- (C);
			\filldraw[fill=purple!80!black, draw=black!80, line width=0.5pt] (C) circle (2.0pt);
			
			% Simbolo di angolo retto in C
			\draw[line width=0.6pt, color=gray!80!black] 
			($(C)!0.22cm!(B)$) -- ($(C)!0.22cm!(B) + ($(C)!0.22cm!(A)-(C)$)$) -- ($(C)!0.22cm!(A)$);
			
			% Arco dell'angolo theta in B
			\draw[line width=0.5pt, color=gray!80!black] (B) +(0, 0.9) arc (90:55:0.9);
			\node[font=\scriptsize\sffamily, color=gray!90!black] at ($(B)+(0.26, 1.25)$) {$\theta$};
			
			% Arco dell'angolo d\theta in A
			\draw[line width=0.6pt, color=orange!90!red] (A) +(-35:1.6cm) arc (-35:{-35-atan2(3.276, 2.294)}:1.6cm);
			
			% Etichette dei lati
			\node[left=4pt, font=\small\bfseries, color=orange!90!red] at ($(A)!0.5!(B)$) {$\mathrm{d}y$};
			\node[below right=-1pt, font=\small\bfseries, color=purple!80!black] at ($(B)!0.5!(C)$) {$r\,\mathrm{d}\theta$};
			
		\end{scope}
		
		% Badge esplicativo superiore
		\node[badge] at (0.8, 3.8) {$r\,\mathrm{d}\theta = \mathrm{d}y \cos\theta \implies \mathrm{d}y = \dfrac{r\,\mathrm{d}\theta}{\cos\theta}$};
		
	\end{tikzpicture}
	\caption{Ingrandimento del triangolo rettangolo infinitesimo in prossimità dell'elemento $\mathrm{d}y$, che evidenzia la componente perpendicolare $r\,\mathrm{d}\theta$.}
	\label{fig:ingrandimento_dy_infografica}
\end{figure}

Si prenda un elemento di linea $dx$ centrato in un punto qualsiasi $P'$. Esiste sempre un secondo elemento, collocato in posizione simmetrica rispetto a $O$, che produce in $P$ un campo dello stesso modulo. Le componenti dei due contributi parallele alla retta sono uguali e opposte e si annullano a coppie, mentre restano le componenti lungo la normale $OP$, che si sommano. Per simmetria, dunque, il campo risultante è perpendicolare alla distribuzione e il suo modulo si ottiene sommando le sole componenti $dE_{oz}$ lungo $OP$.

Indicata con $r$ la distanza tra $P$ e $P'$, e con $\theta$ l'angolo tra $PP'$ e $OP$, la componente utile del campo dovuto a $dx$ vale
\begin{equation}
	dE_{oz}=\frac{1}{4\pi\epsilon_{o}}\,\frac{\cos\theta}{r^{2}}\,dq=\frac{\lambda}{4\pi\epsilon_{o}}\,\frac{\cos\theta}{r^{2}}\,dy .
\end{equation}
Qui $\lambda$ si misura in $\text{C/m}$, $r$ e $dx$ in metri e $dE_{oz}$ in $\text{N/C}$.

Dalla geometria della Fig. 12 si ricavano le relazioni
\begin{equation}
	r=\frac{a}{\cos\theta},\qquad dy=\frac{r\,d\theta}{\cos\theta}=\frac{a\,d\theta}{\cos^{2}\theta},
\end{equation}
con cui l'espressione precedente diventa
\begin{equation}
	dE_{oz}=\frac{\lambda}{4\pi\epsilon_{o}}\,\frac{\cos\theta\cdot\cos^{2}\theta}{a^{2}}\cdot\frac{a\,d\theta}{\cos^{2}\theta}=\frac{\lambda}{4\pi\epsilon_{o}}\,\frac{\cos\theta}{a}\,d\theta .
\end{equation}
Il filo è illimitato, quindi l'angolo $\theta$ varia tra $-\pi/2$ e $\pi/2$. Integrando su tutta la distribuzione si ottiene
\begin{equation}
	E_{o}=\frac{\lambda}{4\pi\epsilon_{o}\,a}\int_{-\pi/2}^{\pi/2}\cos\theta\,d\theta=\frac{\lambda}{4\pi\epsilon_{o}\,a}\cdot 2=\frac{1}{2\pi\epsilon_{o}}\,\frac{\lambda}{a}.
\end{equation}
Il campo decresce come $1/a$ e non dipende dalla posizione lungo il filo; il suo verso è uscente e perpendicolare al filo se $\lambda>0$ ed entrante se $\lambda<0$.


\titoletto{Anello uniformemente carico}
Sull'asse di un anello carico di raggio $R$ con carica totale $Q=2\pi R\lambda$, la distanza da un punto dell'asse vale
\[
r=\sqrt{x^2+R^2}.
\]

\begin{figure}[htbp]
	\centering
	\begin{tikzpicture}[>=Stealth, font=\small\sffamily]
		
		% Stili grafici
		\tikzset{
			gridline/.style={draw=gray!20, line width=0.3pt, dashed},
			badge/.style={draw=gray!40, fill=gray!10, rounded corners=3pt, inner xsep=8pt, inner ysep=4pt, font=\small\sffamily},
			fieldarrow/.style={line width=1.3pt, ->},
			posline/.style={draw=gray!50, line width=0.7pt, dashed},
			axisline/.style={draw=gray!70, line width=0.8pt},
			ringline/.style={draw=orange!80!black, line width=1.5pt},
			dimline/.style={line width=0.6pt, <->, color=gray!80!black}
		}
		
		% Griglia tecnica di sfondo
		\draw[gridline] (-2.2,-3.6) grid (6.8,4.4);
		
		% Coordinate principali
		\coordinate (O)       at (0, 0);          % Centro dell'anello
		\coordinate (P)       at (3.2, 0);        % Punto P sull'asse x
		\coordinate (Pprime)  at (0, 2.4);        % Elemento dl superiore (y = +R)
		\coordinate (Pdouble) at (0, -2.4);       % Elemento dl' inferiore simmetrico (y = -R)
		
		% Assi cartesiani
		\draw[axisline, ->] (-1.5, 0) -- (6.5, 0) node[right=2pt, font=\small\bfseries\itshape, color=gray!80!black] {$x$};
		\draw[axisline, ->] (0, -3.2) -- (0, 3.8) node[above=2pt, font=\small\bfseries\itshape, color=gray!80!black] {$y$};
		
		% Anello in prospettiva 3D (ellisse nel piano yz)
		% Parte posteriore dell'anello (tratteggiata per tridimensionalita)
		\draw[ringline, dashed, color=orange!70!black] (0, 2.4) arc (90:270:0.6cm and 2.4cm);
		% Parte anteriore dell'anello (continua)
		\draw[ringline, color=orange!90!black] (0, -2.4) arc (-90:90:0.6cm and 2.4cm);
		
		% Anelli concentrici di campo attorno agli elementi dl e dl'
		\foreach \r in {0.25, 0.45, 0.65} {
			\draw[orange!50!red, line width=0.4pt, dashed] (Pprime) circle (\r);
			\draw[teal!50, line width=0.4pt, densely dotted] (Pdouble) circle (\r);
		}
		
		% Evidenziazione degli elementi infinitesimi dl e dl' sull'anello
		\draw[line width=3.5pt, color=orange!90!red] (-0.12, 2.37) arc (100:80:0.6cm and 2.4cm);
		\draw[line width=3.5pt, color=teal!70!black] (-0.12, -2.37) arc (-100:-80:0.6cm and 2.4cm);
		
		% Congiungenti r dai punti sorgente al punto P
		\draw[posline] (Pprime) -- (P);
		\draw[posline] (Pdouble) -- (P);
		
		% Etichetta della distanza r
		\node[above right=-1pt, font=\small\itshape, color=gray!80!black] at ($(Pprime)!0.5!(P)$) {$r = \sqrt{x^2 + R^2}$};
		
		% Quota del raggio R sull'asse y
		\draw[dimline] (-0.8, 0) -- (-0.8, 2.4) node[midway, fill=white, inner sep=1.5pt, font=\small\itshape] {$R$};
		\draw[gray!60, line width=0.5pt] (-0.2, 0) -- (-1.0, 0);
		\draw[gray!60, line width=0.5pt] (-0.2, 2.4) -- (-1.0, 2.4);
		
		% Quota della coordinata x sull'asse x
		\draw[dimline] (0, -1.0) -- (3.2, -1.0) node[midway, fill=white, inner sep=1.5pt, font=\small\itshape] {$x$};
		\draw[gray!60, line width=0.5pt] (0, -0.3) -- (0, -1.2);
		\draw[gray!60, line width=0.5pt] (3.2, -0.3) -- (3.2, -1.2);
		
		% Vettori campo infinitesimo applicati in P
		% dE_0 da P' (uscente, diretto verso basso-destra)
		\draw[fieldarrow, color=orange!90!red] (P) -- ($(P)!-1.8cm!(Pprime)$) node[below right=-1pt, font=\bfseries] {$\mathrm{d}\vec{E}_0$};
		% dE_0' dall'elemento simmetrico P'' (uscente, diretto verso alto-destra)
		\draw[fieldarrow, color=teal!80!black] (P) -- ($(P)!-1.8cm!(Pdouble)$) node[above right=-1pt, font=\bfseries] {$\mathrm{d}\vec{E}_0'$};
		
		% Componente risultante lungo l'asse x dE_x
		\coordinate (E1) at ($(P)!-1.8cm!(Pprime)$);
		\coordinate (E2) at ($(P)!-1.8cm!(Pdouble)$);
		\coordinate (Eres) at ($(E1)+(E2)-(P)$);
		\draw[gray!40, line width=0.6pt, dashed] (E1) -- (Eres) -- (E2);
		\draw[fieldarrow, color=purple!80!black] (P) -- (Eres) node[right=2pt, font=\bfseries] {$\mathrm{d}\vec{E}_x$};
		
		% Archi per l'angolo theta
		% 1. Angolo theta al vertice P tra r e l'asse x
		\draw[line width=0.5pt, color=gray!80!black] (P) +(-0.9,0) arc (180:143.13:0.9);
		\node[font=\scriptsize\sffamily, color=gray!90!black] at ($(P)+(-1.15, 0.22)$) {$\theta$};
		
		% 2. Angolo theta tra dE_x e dE_0
		\draw[line width=0.5pt, color=gray!80!black] (P) +(0.9,0) arc (0:-36.87:0.9);
		\node[font=\scriptsize\sffamily, color=gray!90!black] at ($(P)+(1.15, -0.22)$) {$\theta$};
		
		% 3. Angolo theta tra dE_x e dE_0'
		\draw[line width=0.5pt, color=gray!80!black] (P) +(0.9,0) arc (0:36.87:0.9);
		\node[font=\scriptsize\sffamily, color=gray!90!black] at ($(P)+(1.15, 0.22)$) {$\theta$};
		
		% Punti ed etichette
		\filldraw[fill=black!70, draw=black] (O) circle (1.8pt) node[below left=2pt, font=\small\bfseries] {$O$};
		\filldraw[fill=black!70, draw=black] (P) circle (1.8pt) node[above=4pt, font=\small\itshape] {$P$};
		
		\filldraw[fill=orange!80!red, draw=black!80, line width=0.6pt] (Pprime) circle (3.0pt);
		\node[right=4pt, font=\small\bfseries, color=orange!90!red] at (Pprime) {$\mathrm{d}l$};
		\node[left=2pt, font=\small\bfseries] at (Pprime) {$P'$};
		
		\filldraw[fill=teal!70!black, draw=black!80, line width=0.6pt] (Pdouble) circle (3.0pt);
		\node[right=4pt, font=\small\bfseries, color=teal!80!black] at (Pdouble) {$\mathrm{d}l'$};
		
		% Badge descrittivo superiore
		\node[badge] at (3.5, 3.7) {$\mathrm{d}E_x = \mathrm{d}E_0 \cos\theta = \dfrac{1}{4\pi\varepsilon_0} \dfrac{x\,\mathrm{d}q}{(x^2+R^2)^{3/2}}$};
		
	\end{tikzpicture}
	\caption{Calcolo del campo elettrostatico $\mathrm{d}\vec{E}_x$ lungo l'asse $x$ di un anello carico uniformemente di raggio $R$ per sovrapposizione degli elementi simmetrici $\mathrm{d}l$ e $\mathrm{d}l'$.}
	\label{fig:campo_anello_asse_infografica}
\end{figure}

Poichè
\[
\cos\theta =\frac{x}{r} =\frac{x}{(x^2+R^2)^{1/2}} 
\]
La componente assiale infinitesima si scrive
\[
\dd E_x=\frac{1}{4\pi\epz}\,\frac{\lambda\,x\,\dd l}{(x^2+R^2)^{3/2}} 
\]
e l'integrazione produce
\[
E_x=\frac{1}{4\pi\epz}\,\frac{Q\,x}{(x^2+R^2)^{3/2}} .
\]
Per $x\gg R$ si ottiene il campo di una carica puntiforme
\[
E_x\simeq\frac{Q}{4\pi\epz x^2},
\]
mentre per $x\ll R$ il campo diventa lineare nella posizione
\[
E_x\simeq\frac{Q\,x}{4\pi\epz R^3}.
\]
Una carica $q<0$ posta in prossimità del centro subisce la forza di tipo elastico
\[
F_x=qE_x=-\frac{|q|Q}{4\pi\epz R^3}\,x,
\]
dando origine a un moto armonico di pulsazione
\[
\omega=\sqrt{\frac{|q|Q}{4\pi\epz mR^3}}.
\]

\titoletto{Disco uniformemente carico e piano indefinito}
Per il disco di raggio $R$ e carica $q=\sigma\pi R^2$, l'elemento d'area d'anello ha carica
\[
\dd S=2\pi s\,\dd s, \qquad \dd q=\sigma\,2\pi s\,\dd s,
\]

\begin{figure}[htbp]
	\centering
	\begin{tikzpicture}[>=Stealth, font=\small\sffamily]
		
		% Stili grafici
		\tikzset{
			gridline/.style={draw=gray!20, line width=0.3pt, dashed},
			badge/.style={draw=gray!40, fill=gray!10, rounded corners=3pt, inner xsep=8pt, inner ysep=4pt, font=\small\sffamily},
			fieldarrow/.style={line width=1.3pt, ->},
			posline/.style={draw=gray!50, line width=0.7pt, dashed},
			axisline/.style={draw=gray!70, line width=0.8pt},
			diskfill/.style={fill=orange!12, draw=orange!70!black, line width=1.0pt},
			annulusfill/.style={fill=orange!60!red, draw=orange!90!red, line width=0.8pt},
			dimline/.style={line width=0.6pt, <->, color=gray!80!black}
		}
		
		% Griglia tecnica di sfondo
		\draw[gridline] (-2.2,-3.6) grid (6.8,4.4);
		
		% Coordinate principali
		\coordinate (O)       at (0, 0);          % Centro del disco
		\coordinate (P)       at (3.2, 0);        % Punto P sull'asse z
		\coordinate (Pprime)  at (0, 1.5);        % Elemento ds superiore (s = +1.5)
		\coordinate (Pdouble) at (0, -1.5);       % Elemento ds inferiore simmetrico
		
		% Assi cartesiani
		\draw[axisline, ->] (-1.5, 0) -- (6.5, 0) node[right=2pt, font=\small\bfseries\itshape, color=gray!80!black] {$z$};
		\draw[axisline, ->] (0, -3.2) -- (0, 3.8) node[above=2pt, font=\small\bfseries\itshape, color=gray!80!black] {$s$};
		
		% Disegno del disco intero di raggio R in prospettiva 3D (ellisse rx=0.6cm, ry=2.4cm)
		\filldraw[diskfill] (O) ellipse (0.6cm and 2.4cm);
		
		% Corona circolare infinitesima dS di raggio s e spessore ds
		\filldraw[annulusfill, even odd rule] 
		(O) ellipse (0.41cm and 1.62cm)
		(O) ellipse (0.35cm and 1.38cm);
		
		% Tratto posteriore tratteggiato del bordo interno del disco per tridimensionalita
		\draw[gray!60, line width=0.6pt, dashed] (0, 2.4) arc (90:270:0.6cm and 2.4cm);
		
		% Anelli concentrici di campo attorno agli elementi della corona circolare
		\foreach \r in {0.25, 0.45, 0.65} {
			\draw[orange!50!red, line width=0.4pt, dashed] (Pprime) circle (\r);
			\draw[teal!50, line width=0.4pt, densely dotted] (Pdouble) circle (\r);
		}
		
		% Congiungenti r dai punti della corona al punto P
		\draw[posline] (Pprime) -- (P);
		\draw[posline] (Pdouble) -- (P);
		
		% Etichetta della distanza r
		\node[above right=-1pt, font=\small\itshape, color=gray!80!black] at ($(Pprime)!0.55!(P)$) {$r = \sqrt{z^2 + s^2}$};
		
		% Quota del raggio interno s della corona
		\draw[dimline] (-0.8, 0) -- (-0.8, 1.5) node[midway, fill=white, inner sep=1.0pt, font=\small\itshape] {$s$};
		\draw[gray!60, line width=0.5pt] (-0.2, 0) -- (-1.0, 0);
		\draw[gray!60, line width=0.5pt] (-0.2, 1.5) -- (-1.0, 1.5);
		
		% Quota del raggio totale del disco R
		\draw[dimline] (-1.6, 0) -- (-1.6, 2.4) node[midway, fill=white, inner sep=1.0pt, font=\small\itshape] {$R$};
		\draw[gray!60, line width=0.5pt] (-0.2, 2.4) -- (-1.8, 2.4);
		
		% Quota dello spessore ds in cima al disco
		\draw[dimline] (0.2, 1.38) -- (0.2, 1.62);
		\node[right=3pt, font=\scriptsize\bfseries, color=orange!90!red] at (0.2, 1.5) {$\mathrm{d}s$};
		
		% Quota della coordinata z sull'asse z
		\draw[dimline] (0, -1.0) -- (3.2, -1.0) node[midway, fill=white, inner sep=1.5pt, font=\small\itshape] {$z$};
		\draw[gray!60, line width=0.5pt] (0, -0.3) -- (0, -1.2);
		\draw[gray!60, line width=0.5pt] (3.2, -0.3) -- (3.2, -1.2);
		
		% Vettori campo infinitesimo applicati in P
		% dE_0 da P' (uscente, diretto verso basso-destra)
		\draw[fieldarrow, color=orange!90!red] (P) -- ($(P)!-1.8cm!(Pprime)$) node[below right=-1pt, font=\bfseries] {$\mathrm{d}\vec{E}_0$};
		% dE_0' dall'elemento simmetrico P'' (uscente, diretto verso alto-destra)
		\draw[fieldarrow, color=teal!80!black] (P) -- ($(P)!-1.8cm!(Pdouble)$) node[above right=-1pt, font=\bfseries] {$\mathrm{d}\vec{E}_0'$};
		
		% Componente risultante lungo l'asse z dE_z
		\coordinate (E1) at ($(P)!-1.8cm!(Pprime)$);
		\coordinate (E2) at ($(P)!-1.8cm!(Pdouble)$);
		\coordinate (Eres) at ($(E1)+(E2)-(P)$);
		\draw[gray!40, line width=0.6pt, dashed] (E1) -- (Eres) -- (E2);
		\draw[fieldarrow, color=purple!80!black] (P) -- (Eres) node[right=2pt, font=\bfseries] {$\mathrm{d}\vec{E}_z$};
		
		% Archi per l'angolo theta
		% 1. Angolo theta al vertice P tra r e l'asse z
		\draw[line width=0.5pt, color=gray!80!black] (P) +(-0.9,0) arc (180:154.88:0.9);
		\node[font=\scriptsize\sffamily, color=gray!90!black] at ($(P)+(-1.15, 0.18)$) {$\theta$};
		
		% 2. Angolo theta tra dE_z e dE_0
		\draw[line width=0.5pt, color=gray!80!black] (P) +(0.9,0) arc (0:-25.12:0.9);
		\node[font=\scriptsize\sffamily, color=gray!90!black] at ($(P)+(1.15, -0.18)$) {$\theta$};
		
		% 3. Angolo theta tra dE_z e dE_0'
		\draw[line width=0.5pt, color=gray!80!black] (P) +(0.9,0) arc (0:25.12:0.9);
		\node[font=\scriptsize\sffamily, color=gray!90!black] at ($(P)+(1.15, 0.18)$) {$\theta$};
		
		% Punti ed etichette
		\filldraw[fill=black!70, draw=black] (O) circle (1.8pt) node[below left=2pt, font=\small\bfseries] {$O$};
		\filldraw[fill=black!70, draw=black] (P) circle (1.8pt) node[above=4pt, font=\small\itshape] {$P$};
		
		\filldraw[fill=orange!80!red, draw=black!80, line width=0.6pt] (Pprime) circle (3.0pt);
		\node[left=3pt, font=\small\bfseries, color=orange!90!red] at (Pprime) {$\mathrm{d}S$};
		
		\filldraw[fill=teal!70!black, draw=black!80, line width=0.6pt] (Pdouble) circle (3.0pt);
		
		% Badge descrittivo superiore
		\node[badge] at (2.4, 3.7) {$\mathrm{d}E_z = \dfrac{\sigma\,z}{2\varepsilon_0} \dfrac{s\,\mathrm{d}s}{(z^2+s^2)^{3/2}}$};
		
	\end{tikzpicture}
	\caption{Scomposizione in anelli concentrici di area $\mathrm{d}S = 2\pi s\,\mathrm{d}s$ per il calcolo del campo elettrostatico $\mathrm{d}\vec{E}_z$ sull'asse $z$ di un disco carico.}
	\label{fig:campo_disco_infografica}
\end{figure}

generando la componente assiale
\[
\dd E_z=\frac{\sigma\,z}{2\epz}\,\frac{s\,\dd s}{(z^2+s^2)^{3/2}} .
\]
Sfruttando la primitiva
\[
\int\frac{s\,\dd s}{(z^2+s^2)^{3/2}}=-\frac{1}{\sqrt{z^2+s^2}},
\]
si calcola l'integrale definita
\[
E_z=\frac{\sigma}{2\epz}\,\frac{z}{|z|}\left[1-\frac{|z|}{\sqrt{z^2+R^2}}\right].
\]

\begin{figure}[htbp]
	\centering
	\begin{tikzpicture}[>=Stealth, font=\small\sffamily]
		
		% Stili grafici
		\tikzset{
			gridline/.style={draw=gray!20, line width=0.3pt, dashed},
			badge/.style={draw=gray!40, fill=gray!10, rounded corners=3pt, inner xsep=8pt, inner ysep=4pt, font=\small\sffamily},
			axisline/.style={draw=gray!70, line width=0.8pt, ->},
			curveline/.style={draw=purple!85!black, line width=1.6pt},
			refline/.style={draw=gray!60, line width=0.6pt, dashed}
		}
		
		% Griglia tecnica di sfondo
		\draw[gridline] (-5.2,-3.2) grid (5.2,3.8);
		
		% Assi cartesiani
		\draw[axisline] (-4.8, 0) -- (4.8, 0) node[right=2pt, font=\small\bfseries\itshape, color=gray!80!black] {$z$};
		\draw[axisline] (0, -2.8) -- (0, 3.2) node[above=2pt, font=\small\bfseries\itshape, color=gray!80!black] {$E_0(z)$};
		
		% Linee di riferimento per +sigma/2eps0 e -sigma/2eps0
		\draw[refline] (-4.5, 2.0) -- (4.5, 2.0);
		\draw[refline] (-4.5,-2.0) -- (4.5,-2.0);
		
		% Etichette asse y
		\node[left=3pt, font=\small\sffamily] at (0, 2.0) {$+\frac{\sigma}{2\varepsilon_0}$};
		\node[left=3pt, font=\small\sffamily] at (0,-2.0) {$-\frac{\sigma}{2\varepsilon_0}$};
		
		% Grafico del campo del disco
		\draw[curveline, color=orange!90!red, domain=0.08:4.5, samples=80, smooth] 
		plot (\x, {2.0 * (1 - \x / sqrt(\x*\x + 1.2))});
		
		\draw[curveline, color=teal!80!black, domain=-4.5:-0.08, samples=80, smooth] 
		plot (\x, {-2.0 * (1 - abs(\x) / sqrt(\x*\x + 1.2))});
		
		% Punti limite in z = 0
		\filldraw[fill=white, draw=orange!90!red, line width=1.0pt] (0, 2.0) circle (2.2pt);
		\filldraw[fill=white, draw=teal!80!black, line width=1.0pt] (0,-2.0) circle (2.2pt);
		
		% Valore esatto al centro z = 0
		\filldraw[fill=purple!85!black, draw=black] (0, 0) circle (2.2pt) node[below left=2pt, font=\small\bfseries] {$O$};
		
		% Badge esplicativo
		\node[badge] at (3, 2.7) {$E_0(z) = \dfrac{\sigma}{2\varepsilon_0} \dfrac{z}{|z|} \left[1 - \dfrac{|z|}{\sqrt{z^2+R^2}}\right]$};
		
	\end{tikzpicture}
	\caption{Andamento del campo elettrostatico $E_0(z)$ in funzione della distanza $z$ per un disco circolare carico di raggio $R$.}
	\label{fig:erdas_fig48_disco}
\end{figure}

Per $z\to 0^{+}$ e $z\to 0^{-}$ si osservano i limiti
\[
\lim_{z\to 0^{+}} E_z = \frac{\sigma}{2\epz}, \qquad \lim_{z\to 0^{-}} E_z = -\frac{\sigma}{2\epz},
\]
mostrando una discontinuità pari a $\sigma/\epz$. 

Per un piano indefinito ($R\to\infty$) il campo diventa uniforme

\begin{figure}[htbp]
	\centering
	\begin{tikzpicture}[>=Stealth, font=\small\sffamily]
		
		% Stili grafici
		\tikzset{
			gridline/.style={draw=gray!20, line width=0.3pt, dashed},
			badge/.style={draw=gray!40, fill=gray!10, rounded corners=3pt, inner xsep=8pt, inner ysep=4pt, font=\small\sffamily},
			fieldarrow/.style={line width=1.3pt, ->, color=orange!90!red},
			planefill/.style={fill=orange!12, draw=orange!80!black, line width=1.0pt}
		}
		
		% Griglia tecnica di sfondo
		\draw[gridline] (-4.6,-3.6) grid (4.6,3.8);
		
		% Coordinate d'origine delle linee di campo sul piano verticale
		\coordinate (L1) at (-0.3, -2.0);
		\coordinate (L2) at (-0.1, -0.7);
		\coordinate (L3) at (0.1, 0.6);
		\coordinate (L4) at (0.3, 1.9);
		
		% Linee di campo a sinistra (uscenti, x < 0)
		\draw[fieldarrow] (L1) -- ($(L1)+(-2.5,0)$);
		\draw[fieldarrow] (L2) -- ($(L2)+(-2.5,0)$);
		\draw[fieldarrow] (L3) -- ($(L3)+(-2.5,0)$);
		\draw[fieldarrow] (L4) -- ($(L4)+(-2.5,0)$) node[left=2pt, font=\bfseries] {$\vec{E}$};
		
		% Piano indefinito verticale in prospettiva 3D
		\filldraw[planefill] (-0.8, -3.0) -- (0.8, -1.5) -- (0.8, 3.0) -- (-0.8, 1.5) -- cycle;
		
		% Distribuzione di cariche positive (+) sul piano verticale
		\node[color=orange!90!red, font=\bfseries\small] at (-0.2, -2.2) {$+$};
		\node[color=orange!90!red, font=\bfseries\small] at ( 0.3, -1.4) {$+$};
		\node[color=orange!90!red, font=\bfseries\small] at (-0.3, -0.6) {$+$};
		\node[color=orange!90!red, font=\bfseries\small] at ( 0.1,  0.2) {$+$};
		\node[color=orange!90!red, font=\bfseries\small] at (-0.2,  1.1) {$+$};
		\node[color=orange!90!red, font=\bfseries\small] at ( 0.3,  2.0) {$+$};
		
		% Etichetta della densità di carica
		\node[above=2pt, font=\small\bfseries, color=orange!80!black] at (-0.5, 2.5) {$\sigma > 0$};
		
		% Linee di campo a destra (uscenti, x > 0)
		\draw[fieldarrow] (L1) -- ($(L1)+(2.5,0)$);
		\draw[fieldarrow] (L2) -- ($(L2)+(2.5,0)$);
		\draw[fieldarrow] (L3) -- ($(L3)+(2.5,0)$);
		\draw[fieldarrow] (L4) -- ($(L4)+(2.5,0)$) node[right=2pt, font=\bfseries] {$\vec{E}$};
		
		% Badge descrittivo superiore
		\node[badge] at (2, 3.4) {$E = \dfrac{\sigma}{2\varepsilon_0}$};
		
	\end{tikzpicture}
	\caption{Campo elettrostatico uniforme e vettori delle linee di campo $\vec{E}$ generati da un piano indefinito verticale con densità di carica positiva $\sigma$.}
	\label{fig:piano_verticale_infografica}
\end{figure}


\begin{figure}[htbp]
	\centering
	\begin{tikzpicture}[>=Stealth, font=\small\sffamily]
		
		% Stili grafici
		\tikzset{
			gridline/.style={draw=gray!20, line width=0.3pt, dashed},
			badge/.style={draw=gray!40, fill=gray!10, rounded corners=3pt, inner xsep=8pt, inner ysep=4pt, font=\small\sffamily},
			axisline/.style={draw=gray!70, line width=0.8pt, ->},
			curveline/.style={draw=purple!85!black, line width=1.6pt},
			refline/.style={draw=gray!60, line width=0.6pt, dashed}
		}
		
		% Griglia tecnica di sfondo
		\draw[gridline] (-4.8,-3.2) grid (4.8,3.8);
		
		% Assi cartesiani
		\draw[axisline] (-4.4, 0) -- (4.4, 0) node[right=2pt, font=\small\bfseries\itshape, color=gray!80!black] {$z$};
		\draw[axisline] (0, -2.8) -- (0, 3.2) node[above=2pt, font=\small\bfseries\itshape, color=gray!80!black] {$E_0(z)$};
		
		% Linee di riferimento
		\draw[refline] (-4.2, 2.0) -- (4.2, 2.0);
		\draw[refline] (-4.2,-2.0) -- (4.2,-2.0);
		
		% Etichette asse y
		\node[left=3pt, font=\small\sffamily] at (0, 2.0) {$+\frac{\sigma}{2\varepsilon_0}$};
		\node[left=3pt, font=\small\sffamily] at (1.2,-2.0) {$-\frac{\sigma}{2\varepsilon_0}$};
		
		% Tratti orizzontali del campo
		\draw[curveline, color=orange!90!red] (0.08, 2.0) -- (4.2, 2.0);
		\draw[curveline, color=teal!80!black] (-4.2,-2.0) -- (-0.08,-2.0);
		
		% Punti aperti sui limiti
		\filldraw[fill=white, draw=orange!90!red, line width=1.0pt] (0, 2.0) circle (2.2pt);
		\filldraw[fill=white, draw=teal!80!black, line width=1.0pt] (0,-2.0) circle (2.2pt);
		
		% Valore nel piano z = 0
		\filldraw[fill=purple!85!black, draw=black] (0, 0) circle (2.2pt) node[below left=2pt, font=\small\bfseries] {$O$};
		
		% Badge esplicativo
		\node[badge] at (3, 3.2) {$E_0(z) = \begin{cases} +\dfrac{\sigma}{2\varepsilon_0} & \text{per } z > 0 \\[1ex] -\dfrac{\sigma}{2\varepsilon_0} & \text{per } z < 0 \end{cases}$};
		
	\end{tikzpicture}
	\caption{Andamento del campo elettrostatico $E_0(z)$ generato da un piano indefinito uniformemente carico.}
	\label{fig:erdas_fig49_piano}
\end{figure}

\[
E_z = \begin{cases} +\dfrac{\sigma}{2\epz} & \text{per } z>0 \\[1.5ex] -\dfrac{\sigma}{2\epz} & \text{per } z<0. \end{cases}
\]

\titoletto{Doppio strato}
Nel doppio strato con densità $+\sigma$ e $-\sigma$, il campo interno vale

\begin{figure}[htbp]
	\centering
	\begin{tikzpicture}[>=Stealth, font=\small\sffamily]
		
		% Stili grafici
		\tikzset{
			gridline/.style={draw=gray!20, line width=0.3pt, dashed},
			badge/.style={draw=gray!40, fill=gray!10, rounded corners=3pt, inner xsep=8pt, inner ysep=4pt, font=\small\sffamily},
			fieldarrow_pos/.style={line width=1.1pt, ->, color=orange!90!red},
			fieldarrow_neg/.style={line width=1.1pt, ->, color=teal!80!black},
			fieldarrow_tot/.style={line width=1.4pt, ->, color=purple!85!black},
			planefill_pos/.style={fill=orange!12, draw=orange!80!black, line width=1.0pt},
			planefill_neg/.style={fill=teal!10, draw=teal!80!black, line width=1.0pt}
		}
		
		% Griglia tecnica di sfondo
		\draw[gridline] (-5.6,-3.6) grid (5.6,3.8);
		
		% Coordinate d'altezza per le linee di campo nei tre domini (y = -1.2 e y = 0.8)
		% ----------------------------------------------------
		% DOMINIO I: Esterno sinistro (x < -1.8)
		% ----------------------------------------------------
		\coordinate (I1) at (-2.3, 0.8);
		\coordinate (I2) at (-2.3, -1.2);
		
		% E_+ (uscente dal piano positivo, verso sinistra)
		\draw[fieldarrow_pos] (I1) -- ($(I1)+(-1.6,0)$) node[above left=-2pt, font=\bfseries\scriptsize] {$\vec{E}_+$};
		\draw[fieldarrow_pos] (I2) -- ($(I2)+(-1.6,0)$);
		
		% E_- (entrante verso il piano negativo, verso destra)
		\draw[fieldarrow_neg] ($(I1)+(-1.6,0)$) -- (I1) node[midway, below, font=\bfseries\scriptsize] {$\vec{E}_-$};
		\draw[fieldarrow_neg] ($(I2)+(-1.6,0)$) -- (I2);
		
		
		% ----------------------------------------------------
		% DOMINIO II: Regione interna (-1.8 < x < 1.8)
		% ----------------------------------------------------
		\coordinate (II1) at (-1.3, 0.8);
		\coordinate (II2) at (-1.3, -1.2);
		
		% E_+ (verso destra, dal piano positivo al negativo)
		\draw[fieldarrow_pos] (II1) -- ($(II1)+(1.2,0)$) node[above left=-2pt, font=\bfseries\scriptsize] {$\vec{E}_+$};
		\draw[fieldarrow_pos] (II2) -- ($(II2)+(1.2,0)$);
		
		% E_- (verso destra, richiamato dal piano negativo)
		\draw[fieldarrow_neg] ($(II1)+(1.3,0)$) -- ($(II1)+(2.5,0)$) node[above right=-2pt, font=\bfseries\scriptsize] {$\vec{E}_-$};
		\draw[fieldarrow_neg] ($(II2)+(1.3,0)$) -- ($(II2)+(2.5,0)$);
		
		% Campo totale risultante E_tot (prominente, centralizzato)
		\draw[fieldarrow_tot] (-1.3, -0.2) -- (1.3, -0.2) node[midway, above=2pt, font=\bfseries] {$\vec{E}$};
		
		
		% ----------------------------------------------------
		% DOMINIO III: Esterno destro (x > 1.8)
		% ----------------------------------------------------
		\coordinate (III1) at (2.3, 0.8);
		\coordinate (III2) at (2.3, -1.2);
		
		% E_+ (uscente dal piano positivo, verso destra)
		\draw[fieldarrow_pos] (III1) -- ($(III1)+(1.6,0)$) node[above right=-2pt, font=\bfseries\scriptsize] {$\vec{E}_+$};
		\draw[fieldarrow_pos] (III2) -- ($(III2)+(1.6,0)$);
		
		% E_- (entrante verso il piano negativo, verso sinistra)
		\draw[fieldarrow_neg] ($(III1)+(1.6,0)$) -- (III1) node[midway, below, font=\bfseries\scriptsize] {$\vec{E}_-$};
		\draw[fieldarrow_neg] ($(III2)+(1.6,0)$) -- (III2);
		
		
		% ----------------------------------------------------
		% PIANI INDEFINITI VERTICALI IN PROSPETTIVA 3D
		% ----------------------------------------------------
		% Piano positivo (+sigma) a x = -1.8
		\filldraw[planefill_pos] (-2.4, -2.8) -- (-1.2, -1.3) -- (-1.2, 2.7) -- (-2.4, 1.2) -- cycle;
		
		% Piano negativo (-sigma) a x = +1.8
		\filldraw[planefill_neg] (1.2, -2.8) -- (2.4, -1.3) -- (2.4, 2.7) -- (1.2, 1.2) -- cycle;
		
		% Distribuzione di carica sul piano positivo
		\node[color=orange!90!red, font=\bfseries\small] at (-2.0, -1.8) {$+$};
		\node[color=orange!90!red, font=\bfseries\small] at (-1.5, -0.8) {$+$};
		\node[color=orange!90!red, font=\bfseries\small] at (-2.0,  0.2) {$+$};
		\node[color=orange!90!red, font=\bfseries\small] at (-1.5,  1.2) {$+$};
		\node[color=orange!90!red, font=\bfseries\small] at (-2.0,  2.0) {$+$};
		\node[above=2pt, font=\small\bfseries, color=orange!90!red] at (-1.8, 2.2) {$+\sigma$};
		
		% Distribuzione di carica sul piano negativo
		\node[color=teal!80!black, font=\bfseries\small] at (1.6, -1.8) {$-$};
		\node[color=teal!80!black, font=\bfseries\small] at (2.1, -0.8) {$-$};
		\node[color=teal!80!black, font=\bfseries\small] at (1.6,  0.2) {$-$};
		\node[color=teal!80!black, font=\bfseries\small] at (2.1,  1.2) {$-$};
		\node[color=teal!80!black, font=\bfseries\small] at (1.6,  2.0) {$-$};
		\node[above=2pt, font=\small\bfseries, color=teal!80!black] at (1.8, 2.2) {$-\sigma$};
		
		
		% ----------------------------------------------------
		% BADGES ESPLICATIVI NEI TRE DOMINI
		% ----------------------------------------------------
		% Esterno sinistro
		\node[badge] at (-3.8, 3.2) {$E = 0$};
		
		% Regione interna
		\node[badge] at (0, 3.2) {$E = \dfrac{\sigma}{2\varepsilon_0} + \dfrac{\sigma}{2\varepsilon_0} = \dfrac{\sigma}{\varepsilon_0}$};
		
		% Esterno destro
		\node[badge] at (3.8, 3.2) {$E = 0$};
		
	\end{tikzpicture}
	\caption{Campo elettrostatico generato da un doppio strato di carica con densità $+\sigma$ e $-\sigma$: elisione dei campi all'esterno ($E = 0$) e somma concordante all'interno ($E = \sigma/\varepsilon_0$).}
	\label{fig:doppio_strato_infografica}
\end{figure}



\[
E=\frac{\sigma}{2\epz}+\frac{\sigma}{2\epz}=\frac{\sigma}{\epz},
\]
mentre all'esterno si annulla
\[
E=0.
\]

\begin{figure}[htbp]
	\centering
	\begin{tikzpicture}[>=Stealth, font=\small\sffamily]
		
		% Stili grafici
		\tikzset{
			gridline/.style={draw=gray!20, line width=0.3pt, dashed},
			badge/.style={draw=gray!40, fill=gray!10, rounded corners=3pt, inner xsep=8pt, inner ysep=4pt, font=\small\sffamily},
			axisline/.style={draw=gray!70, line width=0.8pt, ->},
			curveline/.style={draw=purple!85!black, line width=1.6pt},
			refline/.style={draw=gray!60, line width=0.6pt, dashed}
		}
		
		% Griglia tecnica di sfondo
		\draw[gridline] (-5.2,-1.6) grid (5.2,4.2);
		
		% Posizione dei due piani
		\coordinate (z1) at (-2.0, 0);   % Piano positivo (+sigma)
		\coordinate (z2) at ( 2.0, 0);   % Piano negativo (-sigma)
		
		% Assi cartesiani
		\draw[axisline] (-4.8, 0) -- (4.8, 0) node[right=2pt, font=\small\bfseries\itshape, color=gray!80!black] {$z$};
		\draw[axisline] (0, -0.8) -- (0, 3.6) node[above=2pt, font=\small\bfseries\itshape, color=gray!80!black] {$E_0(z)$};
		
		% Linee di riferimento
		\draw[refline] (-2.0, 0) -- (-2.0, 3.0);
		\draw[refline] ( 2.0, 0) -- ( 2.0, 3.0);
		\draw[refline] (-4.5, 2.6) -- (4.5, 2.6);
		\draw[refline] (-4.5, 1.3) -- (4.5, 1.3);
		
		% Etichette
		\node[below=4pt, font=\small\bfseries, color=orange!90!red] at (z1) {$+\sigma$};
		\node[below=4pt, font=\small\bfseries, color=teal!80!black] at (z2) {$-\sigma$};
		\node[left=3pt, font=\small\sffamily] at (0, 2.6) {$\frac{\sigma}{\varepsilon_0}$};
		\node[left=3pt, font=\small\sffamily] at (0, 1.3) {$\frac{\sigma}{2\varepsilon_0}$};
		
		% Grafico del campo
		\draw[curveline, color=teal!80!black] (-4.5, 0) -- (-2.08, 0);
		\draw[curveline, color=orange!90!red] (-1.92, 2.6) -- (1.92, 2.6);
		\draw[curveline, color=teal!80!black] (2.08, 0) -- (4.5, 0);
		
		% Punti aperti sui confini
		\filldraw[fill=white, draw=teal!80!black, line width=1.0pt] (-2.0, 0) circle (2.2pt);
		\filldraw[fill=white, draw=orange!90!red, line width=1.0pt] (-2.0, 2.6) circle (2.2pt);
		\filldraw[fill=white, draw=orange!90!red, line width=1.0pt] (2.0, 2.6) circle (2.2pt);
		\filldraw[fill=white, draw=teal!80!black, line width=1.0pt] (2.0, 0) circle (2.2pt);
		
		% Valori sulle armature
		\filldraw[fill=purple!85!black, draw=black] (-2.0, 1.3) circle (2.4pt);
		\filldraw[fill=purple!85!black, draw=black] (2.0, 1.3) circle (2.4pt);
		
		% Origine
		\filldraw[fill=black!70, draw=black] (0, 0) circle (1.8pt) node[below left=2pt, font=\small\bfseries] {$O$};
		
		% Badge esplicativo
		\node[badge] at (0, 5) {$E_{\text{int}} = \dfrac{\sigma}{\varepsilon_0}, \quad E_{\text{est}} = 0, \quad E_{\text{piano}} = \dfrac{\sigma}{2\varepsilon_0}$};
		
	\end{tikzpicture}
	\caption{Andamento grafico del campo elettrostatico $E_0(z)$ in funzione della coordinata assiale per un doppio strato di carica (Figura 50 delle dispense).}
	\label{fig:erdas_fig50_doppio_strato}
\end{figure}

\titoletto{Strato sferico uniformemente carico}
Per uno strato sferico di raggio $R$ con carica $q=4\pi R^2\sigma$, si hanno le relazioni geometriche

\begin{figure}[htbp]
	\centering
	\begin{tikzpicture}[>=Stealth, font=\small\sffamily]
		
		% Stili grafici
		\tikzset{
			gridline/.style={draw=gray!20, line width=0.3pt, dashed},
			badge/.style={draw=gray!40, fill=gray!10, rounded corners=3pt, inner xsep=8pt, inner ysep=4pt, font=\small\sffamily},
			fieldarrow/.style={line width=1.3pt, ->},
			posline/.style={draw=gray!50, line width=0.7pt, dashed},
			axisline/.style={draw=gray!70, line width=0.8pt},
			ringdashed/.style={draw=orange!90!red, line width=0.8pt, dashed},
			dimline/.style={line width=0.6pt, <->, color=gray!80!black}
		}
		
		% Griglia tecnica di sfondo
		\draw[gridline] (-3.2,-3.2) grid (8.2,3.6);
		
		% Coordinate fondamentali
		\coordinate (O) at (0,0);            % Centro della sfera
		\coordinate (P) at (5.0,0);          % Punto P sull'asse
		
		% Vertici della corona sferica Q e Q' (angolo theta = 48 gradi, R = 2.2 cm)
		\coordinate (Q)  at (1.472, 1.635);
		\coordinate (Q2) at (1.293, 1.780);
		\coordinate (Qprime)  at (1.472, -1.635);
		\coordinate (Qprime2) at (1.293, -1.780);
		
		% 1. Sfera 3D con ombreggiatura radiale e contorno netto
		\shade[ball color=orange!25, opacity=0.42] (O) circle (2.2cm);
		\draw[line width=0.9pt, color=black!85] (O) circle (2.2cm);
		
		% 2. Equatore tridimensionale (sezione mediana in prospettiva)
		\draw[gray!50, line width=0.5pt, dashed] (2.2,0) arc (0:180:2.2cm and 0.55cm);
		\draw[gray!70, line width=0.6pt] (-2.2,0) arc (180:360:2.2cm and 0.55cm);
		
		% 3. Corona sferica ad anello (esclusivamente tratteggiata)
		\draw[ringdashed] (Q) arc (90:-90:0.65cm and 1.635cm);
		\draw[ringdashed] (Q) arc (90:270:0.65cm and 1.635cm);
		\draw[ringdashed] (Q2) arc (90:-90:0.68cm and 1.780cm);
		\draw[ringdashed] (Q2) arc (90:270:0.68cm and 1.780cm);
		
		% Anelli concentrici di campo attorno ai punti sorgente
		\foreach \r in {0.22, 0.42} {
			\draw[orange!50!red, line width=0.4pt, dashed] (Q) circle (\r);
			\draw[teal!50, line width=0.4pt, densely dotted] (Qprime) circle (\r);
		}
		
		% Assi di riferimento e congiungenti
		\draw[axisline, dashed, ->] (O) -- (7.8, 0);
		\draw[posline] (O) -- (Q);
		\draw[posline] (O) -- (Q2);
		\draw[posline] (Q) -- (P);
		\draw[posline] (Qprime) -- (P);
		
		% Raggio R
		\draw[posline] (O) -- (-2.068, -0.752) node[midway, below left=-2pt, font=\small\itshape, color=black] {$R$};
		
		% Linea verticale tratteggiata in P
		\draw[axisline, dashed] (5.0, -1.6) -- (5.0, 1.6);
		
		% Vettori campo elettrostatico applicati in P
		\draw[fieldarrow, color=teal!80!black] (P) -- ($(P)!-1.6cm!(Qprime)$) coordinate (E2) node[above right=-1pt, font=\bfseries] {$\mathrm{d}\vec{E}_0'$};
		\draw[fieldarrow, color=orange!90!red] (P) -- ($(P)!-1.6cm!(Q)$) coordinate (E1) node[below right=-1pt, font=\bfseries] {$\mathrm{d}\vec{E}_0$};
		
		% Parallelogramma di scomposizione
		\coordinate (Eres) at ($(E1)+(E2)-(P)$);
		\draw[gray!50, line width=0.6pt, dashed] (E1) -- (Eres) -- (E2);
		\draw[gray!50, line width=0.5pt, densely dotted] (E1) -- ($(P)!(E1)!(5.0, 1.6)$);
		\draw[gray!50, line width=0.5pt, densely dotted] (E2) -- ($(P)!(E2)!(5.0, 1.6)$);
		
		% Archi angolari
		\draw[line width=0.5pt, color=black] (O) +(0.8,0) arc (0:48:0.8);
		\node[font=\small\sffamily] at (1.05, 0.45) {$\theta$};
		
		\draw[line width=0.5pt, color=black] (P) +(-0.9,0) arc (180:155.15:0.9);
		\node[font=\small\sffamily] at (3.85, 0.22) {$\varphi$};
		
		% Etichetta coordinata x
		\node[above right=1pt, font=\small\itshape, color=black] at ($(Q)!0.45!(P)$) {$x$};
		
		% Quota della distanza r
		\draw[dimline] (0.0, -1.6) -- (5.0, -1.6) node[midway, fill=white, inner sep=2pt, font=\small\itshape] {$r$};
		\draw[gray!60, line width=0.5pt] (0.8, -1.2) -- (0.8, -1.9);
		\draw[gray!60, line width=0.5pt] (5.0, -0.3) -- (5.0, -1.9);
		
		% Punti ed etichette
		\filldraw[fill=black] (O) circle (1.5pt) node[below left=1pt, font=\small\bfseries] {$O$};
		\filldraw[fill=black] (P) circle (1.5pt) node[above left=1pt, font=\small\itshape] {$P$};
		\node[above right=1pt, font=\small\bfseries] at (Q) {$Q$};
		\node[below right=1pt, font=\small\bfseries] at (Qprime) {$Q'$};
		
		% Badge esplicativo
		\node[badge] at (2.5, 3.2) {$r^2 = x^2 + R^2 - 2xR\cos\theta$};
		
	\end{tikzpicture}
	\caption{Geometria d'integrazione e scomposizione vettoriale per il calcolo del campo elettrostatico generato da uno strato sferico.}
	\label{fig:strato_sferico_campo_elettrostatico}
\end{figure}

\[
x^2=r^2+R^2-2rR\cos\theta, \qquad R^2=x^2+r^2-2xr\cos\varphi,
\]
con coseno dato da
\[
\cos\varphi=\frac{x^2+r^2-R^2}{2xr}.
\]
Differenziando la prima espressione si ricava
\[
\sin\theta\,\dd\theta=\frac{x\,\dd x}{rR},
\]
e l'elemento di campo diventa
\[
\dd E=\frac{\sigma R}{4\epz r^2}\left[1+\frac{r^2-R^2}{x^2}\right]\dd x .
\]
L'integrazione mediante la primitiva
\[
\int\left[1+\frac{r^2-R^2}{x^2}\right]\dd x = x - \frac{r^2-R^2}{x}
\]
fornisce per i punti esterni ($r>R$)
\[
E=\frac{q}{4\pi\epz x^2},
\]
mentre nei punti interni ($r<R$) si ottiene
\[
E=0.
\]

\begin{richiamo}{Linee di forza del campo}
Le linee di forza sono curve la cui tangente è parallela ed concorde al campo. Si addensano dove il campo è più intenso, non si incrociano mai e partono dalle cariche positive per terminare in quelle negative.\end{richiamo}

\titoletto{Esercizio guida: corona circolare carica}
L'esercizio finale della corona circolare di raggi $a=2{,}0\,\mathrm{cm}$ e $b=4{,}0\,\mathrm{cm}$ produce il campo
\[
E_z=\frac{\sigma\,z}{2\epz}\left[\frac{1}{\sqrt{z^2+a^2}}-\frac{1}{\sqrt{z^2+b^2}}\right].
\]
Espandendo per $|z|\ll a$ mediante la relazione
\[
\frac{1}{\sqrt{z^2+a^2}}\simeq\frac1a\left(1-\frac{z^2}{2a^2}\right),
\]
si approssima la forza richiamante
\[
F_z=qE_z=-k\,z,
\]
con costante
\[
k=\frac{|q|\,\sigma}{2\epz}\left(\frac1a-\frac1b\right)\simeq 1,4\cdot 10^{-3}\,\mathrm{N/m}.
\]
L'equazione del moto $m\ddot z=-kz$ descrive un oscillatore armonico di pulsazione e periodo
\[
\omega=\sqrt{\frac{k}{m}}\simeq 38\,\mathrm{rad/s},\qquad
T=\frac{2\pi}{\omega}\simeq 0,17\,\mathrm{s}.
\]

\end{document}